\documentclass[11pt,a4paper]{article}
\usepackage[utf8]{inputenc}
\usepackage[T1]{fontenc}
\usepackage{lmodern,amsmath,amssymb,amsthm,mathtools,booktabs,enumitem,microtype}
\usepackage[margin=1in]{geometry}
\usepackage{natbib}
\usepackage{xr-hyper}
\usepackage{hyperref}
\usepackage{fancyhdr,listings}
\hypersetup{pdftitle={Relational Information Accounting},
 pdfauthor={Matthew Parslow},
 pdfsubject={Whole-system informational consequences, ethics, and moral injury},
 pdfkeywords={information ethics, reachability, empowerment, compression}}
\externaldocument[GE-]{01_law_of_graphic_equalisation_theory_referenced}
\externaldocument[MI-]{01b_graphic_equalisation_mind_and_self_grounding}
\externaldocument[SM-]{02_seldon_manifolds_theory_referenced}
\externaldocument[NA-]{04_narrativium_theory_referenced}
\externaldocument[PS-]{02b_seldon_manifolds_social_and_psychological_application}
\newtheorem{definition}{Definition}[section]
\newtheorem{proposition}[definition]{Proposition}
\newtheorem{hypothesis}[definition]{Hypothesis}
\theoremstyle{remark}
\newtheorem{remark}[definition]{Remark}
\newcommand{\R}{\mathcal R}
\newcommand{\C}{\mathcal C}
\newcommand{\I}{\mathcal I}
\newcommand{\Lost}{\operatorname{Lost}}
\newcommand{\Gained}{\operatorname{Gained}}
\newcommand{\NULL}{\mathsf{undefined}}
\DeclareMathOperator{\supp}{supp}
\DeclareMathOperator*{\argmin}{arg\,min}
\lstset{language=Python,basicstyle=\ttfamily\footnotesize,breaklines=true,
 columns=fullflexible,keepspaces=true,showstringspaces=false,frame=single}
\title{Relational Information Accounting\\[5pt]
\large Ethics as an accounting frame for whole-system consequences and reachable futures}
\author{Matthew Parslow\\Independent Researcher}
\date{\textbf{Draft \texorpdfstring{\textemdash}{---}\ \PaperDraft}\\[2pt]\PaperDisplayDate\ \texorpdfstring{\textbullet}{\textperiodcentered}\ document revision \PaperRevision}
\pagestyle{fancy}
\fancyhf{}
\fancyhead[L]{\small Relational Information Accounting}
\fancyfoot[L]{\small\textsc{\PaperDraft\,r\PaperRevision}}
\fancyfoot[R]{\thepage}
\setlength{\headheight}{14pt}
% GENERATED by tools/build_paper_metadata.py from paper_metadata.json.
% DO NOT EDIT: edit paper_metadata.json, then regenerate this file.
\providecommand{\PaperDraft}{draft-2026-09-06.1}
\providecommand{\PaperDisplayDate}{6 September 2026}
\expandafter\def\csname PaperRevision@00_is\endcsname{129}
\expandafter\def\csname PaperRevision@01_law_of_graphic_equalisation_theory_referenced\endcsname{294}
\expandafter\def\csname PaperRevision@01a_graphic_equalisation_physics_grounding\endcsname{47}
\expandafter\def\csname PaperRevision@01b_graphic_equalisation_mind_and_self_grounding\endcsname{47}
\expandafter\def\csname PaperRevision@01c_graphic_equalisation_cosmology_grounding\endcsname{11}
\expandafter\def\csname PaperRevision@01d_graphic_equalisation_born_recovery\endcsname{6}
\expandafter\def\csname PaperRevision@01e_graphic_equalisation_communication_grounding\endcsname{6}
\expandafter\def\csname PaperRevision@01f_graphic_equalisation_qualia_and_consciousness_grounding\endcsname{3}
\expandafter\def\csname PaperRevision@01g_graphic_equalisation_total_perspective_vortex_grounding\endcsname{4}
\expandafter\def\csname PaperRevision@01h_graphic_equalisation_information_gravity_quantum_reconciliation\endcsname{5}
\expandafter\def\csname PaperRevision@01i_graphic_equalisation_passive_continuation_momentum_gravity\endcsname{12}
\expandafter\def\csname PaperRevision@01s_graphic_equalisation_statement\endcsname{294}
\expandafter\def\csname PaperRevision@02_seldon_manifolds_theory_referenced\endcsname{31}
\expandafter\def\csname PaperRevision@02a_seldon_manifolds_psychological_grounding\endcsname{15}
\expandafter\def\csname PaperRevision@02b_seldon_manifolds_social_and_psychological_application\endcsname{7}
\expandafter\def\csname PaperRevision@02c_seldon_manifolds_decision_making\endcsname{11}
\expandafter\def\csname PaperRevision@03_law_of_narrative_momentum_theory_referenced\endcsname{32}
\expandafter\def\csname PaperRevision@04_narrativium_theory_referenced\endcsname{23}
\expandafter\def\csname PaperRevision@05_frames_relation_and_the_empty_limit\endcsname{11}
\expandafter\def\csname PaperRevision@06_assessing_an_account_of_this_kind\endcsname{29}
\expandafter\def\csname PaperRevision@07_causal_attitude_networks_on_the_substrate\endcsname{4}
\expandafter\def\csname PaperRevision@08_certified_answers_as_an_audit\endcsname{4}
\expandafter\def\csname PaperRevision@09_adaptability\endcsname{94}
\expandafter\def\csname PaperRevision@09a_adaptability_implementation\endcsname{7}
\expandafter\def\csname PaperRevision@10_applied_graphic_equalisation\endcsname{7}
\expandafter\def\csname PaperRevision@10a_red_queen_race_civilisational_rate_of_change_llms_biological_timescales_anthropocene_transition\endcsname{13}
\expandafter\def\csname PaperRevision@11_relational_information_ethics\endcsname{16}
\expandafter\def\csname PaperRevision@A1_constructive_graphic_equalisation_implementation_referenced\endcsname{8}
\expandafter\def\csname PaperGenerated@01s_graphic_equalisation_statement\endcsname{1}
\edef\PaperMetadataStem{\jobname}
\expandafter\ifx\csname PaperRevision@\PaperMetadataStem\endcsname\relax
  \PackageError{paper_metadata}{No metadata for \PaperMetadataStem}{Add the paper stem to paper_metadata.json.}
  \def\PaperRevision{?}
\else
  \expandafter\let\expandafter\PaperRevision\csname PaperRevision@\PaperMetadataStem\endcsname
\fi
\newcommand{\PaperDateLine}{\textbf{Draft \texorpdfstring{\textemdash}{---}\ \PaperDraft}\\[2pt]\PaperDisplayDate\ \texorpdfstring{\textbullet}{\textperiodcentered}\ document revision \PaperRevision}
\newcommand{\PaperStatementDateLine}{\textbf{Draft \texorpdfstring{\textemdash}{---}\ \PaperDraft}\\[2pt]\PaperDisplayDate\ \texorpdfstring{\textbullet}{\textperiodcentered}\ statement}
\expandafter\ifx\csname PaperGenerated@\PaperMetadataStem\endcsname\relax
  \newcommand{\PaperFooterStamp}{\textsc{\PaperDraft\,r\PaperRevision}}
\else
  \newcommand{\PaperFooterStamp}{\textsc{\PaperDraft\ / statement / from long r\PaperRevision}}
\fi
\newcommand{\PaperNoticeStamp}{\textbf{DRAFT --- \PaperDraft\ r\PaperRevision} (\PaperDisplayDate).}

\begin{document}
\maketitle
\begin{center}
\fbox{\begin{minipage}{0.92\textwidth}\small
\textbf{DRAFT --- \PaperDraft\ r\PaperRevision} (\PaperDisplayDate).
Circulated for comment. The structural results are conditional on the stated
comparison and grounding assumptions. The biological compression interpretation
is a hypothesis.
\end{minipage}}
\end{center}

\begin{abstract}
Actions change the reachable futures and enabling conditions of the systems
they affect. This paper derives a calculable whole-system account of those
informational consequences. Ethics becomes an accounting frame for consequences
that narrower descriptions of an action leave outside their account. An action removes or introduces reachable
capabilities, changes the resources required for continuation, and imposes
effects along its course. These consequences occur in affected relations
whether anyone recognises them or assigns them value. Measurement describes
them; decision weights determine how a system responds to them.

Starting from relational existence, admissible transformation, consequential
distinction, and grounded expenditure, the paper derives separate loss, gain,
resource, and path accounts. No additional biological, sentience, preference,
or observer-recognition requirement determines eligibility. The construction
retains affected identities through termination and does not assume a common
unit across agents. It provides finite calculation functions, proves
independence from passive evaluative observation and from decision weights,
and recovers a channel-capacity expression for deterministic controllable
futures. Worked cases include protective restriction, restoration, and the
contrast between disposable game encounters and persistent companions. The
account also includes consequences borne by actors and witnesses, providing a
structural interpretation of moral injury through damaged self-relations, trust,
and continuing participation. It also quantifies a conditional implied decision
conversion and the separate consequences of making that conversion explicit.

A further hypothesis interprets human and animal ethical behaviour as compressed
expressions of this informational principle. Exact and approximate conditions
for preserving decisions under compression are given without claiming to recover
a neural implementation. The contribution is an implementable consequence
account and its derivation, situated against information ethics, empowerment,
and reachability-based impact measures; it does not impose a universal ordering
of the resulting decisions.
\end{abstract}

\section{Consequences before decisions}\label{sec:introduction}

A protective intervention can restrict one person's freedom of action while
restoring another's. Both effects occur. Whether the intervention is justified
depends on how those effects are weighed. A theory that records only the chosen
justification loses the restriction; a theory that records only the restriction
loses the protection. A whole-system account must hold both, together with
consequences for the actor, witnesses, and the relations that enable later
coordination. Ethics brings otherwise omitted consequences into view and studies how agents
respond to them through their relative priorities.

The foundational proposition of this paper is:
\begin{quote}
An action changes the reachable futures and enabling conditions of affected
systems; those changes are accountable across the relations in which they occur.
\end{quote}
Here \emph{reachable futures} includes the conditions, resources, and continuing
relations through which a future is available. Destroying a capability, making
its exercise more costly, and temporarily closing it are different changes.
They require different entries even when the final state looks the same.

The resulting lens separates three things. A consequence \emph{occurs} through
an interaction. A measurement \emph{represents} that consequence under explicit
comparison conditions. A decision \emph{evaluates} the account using weights,
priorities, or constraints. An unobserved loss remains a loss. Assigning it zero
decision weight does not restore what was removed.

\begin{center}
\begin{tabular}{@{}p{0.17\textwidth}p{0.43\textwidth}p{0.31\textwidth}@{}}
\toprule
Layer & Determining relation & Role of evaluation \\
\midrule
Occurrence & Actual transformation of capabilities and resource conditions & No recognition required \\
Measurement & Baseline, boundary, identity, horizon, distinction and unit & Specifies the account being taken \\
Decision & Rule applied to the measured or estimated account & Supplies relative weights and priorities \\
\bottomrule
\end{tabular}
\end{center}

The proposal locates informational cost in the first layer. The second supplies
functions for calculating it. The third allows differences in ethical decisions
without changing the consequences already accounted for. An actor's preferences
can causally change what the actor does; once the compared actions are fixed,
changing an evaluator's preferences changes their ranking rather than their
consequence records.

The foundational object is whole-system accounting. The proposed reason for an
ethical frame is an accounting deficit: a narrower description holds the actor's
objective, immediate output, or direct expenditure while leaving other
consequences outside its account. The omitted effects appear as \emph{extra}
ethical considerations because the narrower frame has not included their
bearers, relations, or continuing costs.

Under the relational ontology, those effects are already changes in the system.
A recipient's lost future, an actor's damaged self-relations, a witness's altered
trust, and an institution's impaired coordination all have a place in the same
informational machinery. They require no separate ethical substance or an
observer's assignment of value to become consequences. Ethics therefore becomes
an \emph{accounting frame}: it directs attention to effects and incidence that
an incomplete operational account leaves out. Relative weights determine how
an agent responds to the resulting account.

This is the paper's structural explanation of the need to study ethics. Ethical
language and practice can make hidden costs legible and compress sensitivity to
them. Studying those practices exposes which consequences a system tracks,
which it omits, and how its priorities govern a response. The claim concerns the
function of the ethical frame in this ontology; a chronology of the emergence
of moral language is a separate empirical question.

Here \emph{whole-system} means seeking causal coverage across affected relations,
including the acting system and its internal and social organisation. A finite
account still declares its boundary and unresolved effects. The phrase does not
certify that every remote consequence has been discovered, and it does not
license a single sum over incompatible units or overlapping bearers.

An ethical justification cannot make an incurred cost disappear. An action can
be chosen because it protects others while still damaging its actor's trust,
self-relations, or continuing participation. Moral injury is therefore a
substantive case for the framework: costs can arise within the system that acts,
in a witness, and in the institutions that connect them. The relevant injury
is a change of organisation and continuation, beyond an external judgement that
an action was good or bad.

Three contributions follow. First, the paper derives an attributed cost account
without an observer conferring value on its bearers, including actors and
witnesses as affected systems. Second, it supplies actual
finite functions for computing the account. Third, it proposes that ethical
behaviour can express the underlying principle through compression, without
conscious ethical reasoning or literal execution of these functions. The first
two stand independently of recovery of a biological mechanism.

\section{What the informational foundations supply}\label{sec:foundations}

The derivation needs structures that can be affected, transformations that can
affect them, and distinctions that specify what changes. The preceding papers
supply these objects. This section states the dependencies so that an apparent
derivation cannot conceal a moral premise in its inputs.

\subsection{Existence, representation, and affected frames}

Paper 00 treats existence as standing in relation
\citep{parslow00relation}. A biological organism, a
computation, a simulation, and an imagined structure differ in the relations
they instantiate and the changes they can sustain. The framework gives them
no rank in degree of reality. Being a representation describes a relation
between structures; it does not make the representing structure unreal.

Consequences nevertheless belong to the relations actually instantiated. An
imagined injury instantiates changes in the imagining system and its
representations. It does not thereby instantiate every bodily consequence of
the injury represented. A simulated agent can possess persistent capabilities
and relationships within its environment; a painted depiction can lack them.
Their accounts differ because their continuation structures differ.

An \emph{affected frame} is a structure whose capabilities, expenditure, or
continuing relations change under the comparison. ``Agent'' is used broadly for
such a continuing organisation, without requiring a human body, linguistic
self-report, or deliberative preference. Where agency or individuality is
uncertain, record the candidate partition and its alternatives. Temporal
information provides approaches to individuality across scales
\citep{krakauer2020individuality}; it does not supply a moral rank.

The affected frame is also distinct from an observer taking an account. A tree,
an unmonitored program, and a collective can each be affected without a further
system recognising that effect. If a dimension is instantiated in the affected
relation, an observer is unnecessary for its occurrence. If the dimension is
absent, a label applied by an observer cannot create it.

\subsection{Transformation, capability, and expenditure}

Graphic Equalisation supplies admissible transitions and grounding-dependent
costs \citep{parslow01equalisation}. Paper 01's well-formedness requirement
(Definition~\ref{GE-def:well-formed}) requires quantities to have a frame and
comparison conditions. Its expenditure account
(Proposition~\ref{GE-prop:economics}) concerns changes and their implementation;
it does not identify every expenditure with a decrease of capability.

Seldon Manifolds supplies sets of reachable continuations
\citep{parslow02manifolds}. A goal-filtered
acceptable manifold (Paper 02, Definition~\ref{SM-def:acceptable}) is a further
construction. The present cost account uses reachable continuations before
filtering them by a preferred goal. A frame with no represented preference can
still lose a continuation.

Narrativium supplies a measure of consequentially reachable deformation
\citep{parslow04narrativium}
(Paper 04, Definitions~\ref{NA-def:reachable-deformation}
and~\ref{NA-def:narrativium}). Its distinction between measurement and
expenditure is retained. Its restriction on sums across frames
(Proposition~\ref{NA-prop:no-sum}) prevents adding unlike quantities merely
because both have been called costs.

These are inherited informational dependencies. The construction below adds
accounting definitions and a particular finite realisation. It does not infer
an ethical prohibition from a positive expenditure, nor derive all the
underlying relational machinery from Shannon's communication theory alone.

\subsection{The derivation boundary}

The argument has the following dependency order:
\begin{enumerate}[leftmargin=*]
\item An instantiated relational organisation has admissible continuations.
\item An action changes the transition relation or the state from which those
continuations are available, including the resources required to obtain them.
\item Comparing action and baseline identifies removed and introduced
continuations, and changed expenditure, for each affected relation.
\item These differences constitute the informational consequence account.
\item Ethics is an accounting frame that brings omitted effects and their
bearers into the comparison. Agents respond through represented information,
compressed dispositions, and relative priorities. A decision rule evaluates
the account; its judgement does not constitute the recorded changes.
\end{enumerate}
Steps 1--4 supply the whole-system construction. Step 5 locates ethics as an
accounting frame within that construction and introduces the hypothesis
connecting behaviour with the account.
Nothing in this chain requires a morally approved species, a sentience threshold,
a represented preference, or an observer's award of value. Grounding an actual
quantity is still necessary; an unknown quantity is not a known zero.

\section{The comparison contract}\label{sec:contract}

Cost relative to an action requires a specified alternative. Removing a bridge
has different consequences compared with leaving it open, replacing it, or
closing it before a collapse. A baseline answers which comparison is being
made. It does not decree which comparison must govern every decision.

\begin{definition}[Comparison contract]\label{def:contract}
A contract $\theta$ specifies:
\begin{enumerate}[leftmargin=*]
\item the compared action $a$ and baseline $b$, including subsequent policies;
\item the causal boundary and affected identities $\I_\theta$, with continuity
conditions across both branches;
\item a horizon and observation times or intervals;
\item each frame's consequential outcome space $\Omega_i$, an aligned outcome
map, and a nonnegative measure $\mu_i$ where defined;
\item resource dimensions, units, payers and attribution rules;
\item model uncertainty and conditions under which the account is complete.
\end{enumerate}
\end{definition}

A contract represents the relations being compared. It does not bring their
changes into existence. Moving a measurement boundary changes what an account
contains; it does not undo a consequence outside that boundary. Claims of
completeness therefore require a causal argument, rather than merely a short
list of preferred recipients.

\subsection{Consequential distinctions}

Two descriptions count as the same outcome when the declared grounding identifies
them with the same relevant continuing capability. Renaming an option cannot
create an additional option. Conversely, two equally numerous outcome sets need
not offer the same capabilities. A chance to leave, a chance to work, and a
chance to obtain food are not interchangeable because each set has three members.

The outcome quotient specifies the resolution of an account. Its justification
must be inspectable in the affected dynamics. A numerical outcome mass can express
a grounded amount of an available capability or a sampling convention. A mass
chosen because a decision-maker likes one outcome more than another is an
evaluation weight and belongs in the decision rule. Calling it $\mu_i$ would
hide valuation inside the calculation.

A measure can therefore be objective within its declared grounding without
being the unique measure of every consequence. The full set changes retain
identities that a scalar summary loses. A coarse account that merges a relevant
distinction is incomplete at that finer resolution; its zero does not establish
that no cost occurred.

\subsection{Identity, time, and attribution}

Use the union of affected identities across branches. The loss of a frame cannot
remove that frame from the population being compared. A created successor is
matched to the baseline frame only where the continuity conditions establish
that match. Similar appearance, a copied name, or a new numerical identifier
is insufficient by itself.

For shared resources, separate physical expenditure from its attribution. A
single shared bill can be represented once with links to all affected frames.
Allocating that bill among payers requires either grounded incidence or an
explicit allocation convention. Summing the full bill once for every claimant
manufactures expenditure. Collective and member accounts can coexist, but their
entries do not become disjoint by being listed separately.


The affected population includes the actor whenever its action changes its own
organisation or continuations. It also includes witnesses, dependants, and
institutions where the action changes their enabling relations. Direct bodily
impact does not exhaust causal incidence. A memory, expectation, permission, or
trust relation can change later admissibility and expenditure, and therefore
belongs in the model where that effect is grounded. A represented judgement
inside an affected system is a causal state; it is different from the passive
external evaluation excluded in Proposition~\ref{prop:observer}.

Horizon and path are likewise distinct. A restored system can have the same
terminal capabilities while having undergone an interval of restriction.
The comparison must state whether it takes a terminal account, an account along
a course, or both. No undisclosed preference for the present over the future is
introduced by the time index.

\section{Deriving the cost functions}\label{sec:derivation}

The account starts with what remains possible. Expenditure and time enter as
additional typed dimensions because neither is recoverable from a final option
count alone.

\subsection{Reachable futures and their changes}

Let $S_x$ be the state space for branch $x\in\{a,b\}$ and let $T_x$ be its
admissible transition relation. Let $s_x$ be the post-intervention starting state.
For frame $i$, the map $\pi_i$ assigns aligned consequential outcomes to states;
$\bot$ means that the state instantiates no continuing outcome for that frame.
For a fixed-step horizon $h$ define
\begin{equation}\label{eq:reachable}
\R_i(x;h)=\{\pi_i(s):s_x\xrightarrow[T_x]{h}s,\ \pi_i(s)\ne\bot\}.
\end{equation}
A horizon of \emph{up to} $h$ steps uses the union over $0\leq t\leq h$ instead.
These answer different questions. Terminal persistence must be encoded if an
exact-step account is intended to include terminated world states.

Resources that constrain feasibility belong in the state or admissibility
relation: for example $(s,\beta,t)$ can carry remaining budget $\beta$ and time
$t$. Otherwise a traversal can count an option that its bearer cannot afford.
An institution's permission and a system's physical capacity can likewise be
separate state components when either changes the actual reachable set.

\begin{definition}[Lost and gained futures]\label{def:loss}
For aligned outcome spaces under one contract, define
\begin{align}
\Lost_i(a\mid b)&=\R_i(b)\setminus\R_i(a),&
L_i(a\mid b)&=\mu_i(\Lost_i(a\mid b)),\label{eq:loss}\\
\Gained_i(a\mid b)&=\R_i(a)\setminus\R_i(b),&
G_i(a\mid b)&=\mu_i(\Gained_i(a\mid b)).\label{eq:gain}
\end{align}
If $\mu_i$ is unavailable, retain the sets and mark their numerical measures
undefined. For finite distinct outcomes, counting measure is a calculable
specialisation.
\end{definition}

The functions follow directly from the two reachable sets. They require no
preference ordering on their members. The pair $(L_i,G_i)$ is retained rather
than cancelled: an action can close an exit and open an unrelated opportunity.
When finite measures exist,
\begin{equation}\label{eq:net}
\mu_i(\R_i(a))-\mu_i(\R_i(b))=G_i-L_i.
\end{equation}
The difference is a net capability summary. It omits which futures were replaced.

\subsection{The information-theoretic quantity}

Counting consequentially distinct futures connects the construction to
information theory. In a deterministic control channel, let $U$ denote legal
inputs and $Y=f(U)$ their consequential outputs. Empowerment is the channel
capacity
\begin{equation}\label{eq:empowerment}
E=\max_{p(U)}I(U;Y),
\end{equation}
introduced as an agent-centred measure of potential control by
\citet{klyubin2005empowerment}. Its units are bits when logarithms have base two.

\begin{proposition}[Finite deterministic capacity]\label{prop:capacity}
If a deterministic channel has $m\geq1$ distinct outputs and each is attainable
by a legal input, then $E=\log_2m$.
\end{proposition}
\begin{proof}
Determinism gives $H(Y\mid U)=0$, so $I(U;Y)=H(Y)\leq\log_2m$. Choose one input
for each output and assign these $m$ inputs equal probability. The outputs are
then uniform and the upper bound is attained.
\end{proof}

Thus a nested restriction from $m_b$ to $m_a\geq1$ controllable outcomes removes
\begin{equation}\label{eq:bitloss}
\Delta E=\log_2m_b-\log_2m_a=\log_2(m_b/m_a)
\end{equation}
bits of attainable distinction. This recovers a known capacity expression; it
makes the informational content of the finite loss concrete. For non-nested
sets, capacity difference alone conceals replacements. For stochastic systems,
possible outcomes and controllable outcomes differ: an uncontrollable lottery
can have many reachable results without high channel capacity. Use the actual
channel when measuring empowerment, and reachability when measuring possibility.

No channel with zero attainable outputs is assigned $\log_2 0$. Extinction is
recorded through lost continuations and identity, while its capacity entry is
undefined. A deterministic one-output channel has zero bits of control; removing
its sole continuing outcome can still produce a nonzero set loss. Capacity is
one dimension of the account, rather than a complete measure of existence.

Shannon uncertainty alone is also insufficient. A string of random bits does
not acquire ethical priority through its entropy. Relevant information here is
the organisation of consequential continuations and their enabling relations.
Communication theory supplies the capacity calculation
\citep{shannon1948communication}; the relational grounding supplies what the
channel connects and whom its alteration affects.

\subsection{Resource burdens and access prices}

Let $q_{ir}(x)$ be branch $x$'s actual or predicted expenditure borne by frame
$i$ in grounded resource dimension $r$. The incremental burden is
\begin{equation}\label{eq:resource}
K_{ir}(a\mid b)=q_{ir}(a)-q_{ir}(b).
\end{equation}
This is signed: a negative entry records saved expenditure. Energy, time,
attention, computational work, and access payments require their own instantiated
units and incidence. A grounding that supports joules does not automatically
support a quantitative measure of attention.

An \emph{access price} is different. If $p_{ir}(y\mid x)$ is the least expenditure
required for frame $i$ to attain outcome $y$ from branch $x$, a change in that
price can burden the frame before it attempts the outcome. Retain the price
profile over aligned outcomes where it is defined. A feasible set can remain
unchanged while all its members become more expensive. Expenditure in
Eq.~\eqref{eq:resource} is what was paid; an access-price change is what would
have to be paid. The two must not be silently substituted.

The substrate does not imply a universal exchange rate between these resource
entries and lost capabilities. Its information-erasure discussion also does not
license charging a fixed Landauer cost for every logical or relational change.
Logical irreversibility and physical implementation conditions must be supplied
before such a bound is used.

\subsection{Effects along the course and restoration}

For aligned comparison times $t$, calculate
\begin{equation}\label{eq:path}
L_i(t;a\mid b)=\mu_i\bigl(\R_i(b;t)\setminus\R_i(a;t)\bigr).
\end{equation}
Here $\R_i(x;t)$ describes capabilities from the state reached at $t$, under the
same declared continuation horizon. These sets are conditional on a realised or
specified course; uncertain courses require a profile or distribution of
accounts. They are not automatically the unconditional terminal set in
Eq.~\eqref{eq:reachable}.

The matching gain profile is
\begin{equation}\label{eq:path-gain}
G_i(t;a\mid b)=\mu_i\bigl(\R_i(a;t)\setminus\R_i(b;t)\bigr).
\end{equation}
A capability made available only temporarily can disappear from a terminal
snapshot while remaining a gain along the course. Retain both profiles without
cancelling unrelated losses and gains.

Retain loss and gain profiles as the primary path account. Where a clock and an
integration measure $\nu$ are grounded, define exposure
\begin{equation}\label{eq:exposure}
D_i(a\mid b)=\int L_i(t;a\mid b)\,d\nu(t).
\end{equation}
For piecewise constant intervals of duration $\Delta t_k$, this is
$\sum_k\Delta t_k L_i(t_k;a\mid b)$. Its units are capability units multiplied
by clock units. It is not another terminal loss in the same units.
Similarly, where the same clock is grounded, gain exposure is
$D_i^+(a\mid b)=\int G_i(t;a\mid b)\,d\nu(t)$. It has the corresponding
capability-time units and remains separate from loss exposure.

Exact restoration at $T$ can give $L_i(T)=0$ while $D_i>0$ and expenditure is
positive. Paper 02b's deprivation-and-restoration example
(Remark~\ref{PS-rem:ethics}) motivates this distinction. An imposed arousal
variable is not sufficient evidence of experienced suffering; the structural
point is that a terminal snapshot can omit an intervening restriction.

This historical account does not posit an invisible historical force acting
on a restored state. It compares the courses that occurred. It also does not
require that a recipient retain a memory of the restriction. An external record
makes the account recoverable; it does not constitute the past transformation.

\subsection{Repair feasibility}

Define a restoration condition $B_i$ on continuing capabilities and relations.
Let $\mathcal P_i$ be the admissible repair paths from the post-action state to a
state satisfying $B_i$. For one grounded resource unit, define
\begin{equation}\label{eq:repair}
\rho_i(a\mid b)=\inf_{p\in\mathcal P_i}q_i(p),\qquad
\inf\varnothing=+\infty.
\end{equation}
A finite nonnegative-cost graph permits a minimum-cost path calculation. With
multiple incompatible units, retain the feasible resource vectors or their
Pareto frontier; a scalar shortest path requires a declared conversion.

A repair estimate is a prospective liability. It is not expenditure already
incurred, and restoring a terminal capability does not reverse prior exposure.
If repair is carried out, its actual expenditure belongs to the repaired
branch's resource record. Adding the estimate again would charge the same
repair twice.

\subsection{The complete account and its finite realisation}

\begin{definition}[Informational consequence account]\label{def:account}
For a contract $\theta$, the account is the typed family
\begin{equation}\label{eq:account}
\C_\theta(a\mid b)=
\left\{\begin{aligned}
 &\bigl(\Lost_i,\Gained_i,L_i,G_i,(K_{ir})_r,(L_i(t))_t,(G_i(t))_t,\\
 &\qquad(p_{ir}(\cdot\mid a),p_{ir}(\cdot\mid b))_r,\rho_i\bigr)
 \end{aligned}:i\in\I_\theta\right\},
\end{equation}
together with identity, shared incidence, grounded stochastic-transition or
control-channel records where applicable, and uncertainty records.
Support sets do not determine outcome probabilities or control reliability:
changes in those records require their own aligned comparison and are not
represented by a zero set difference. A component
is included only where its dimension is instantiated; numerical absence is
marked undefined. Informational cost denotes the losses and burdens in this
account; gains and savings are retained alongside them. An ethical accounting frame
selects and examines consequences and incidence; decision rules evaluate the
account without redefining the occurrence of its entries. A frame that leaves
a consequence out does not make that consequence absent from the system.
\end{definition}

This is a family of actual functions, not a single untyped score. Appendix
\ref{app:functions} prints a finite Python realisation. Its graph traversal
computes Eq.~\eqref{eq:reachable}; outcome comparison computes
Eqs.~\eqref{eq:loss}--\eqref{eq:gain}; resource comparison computes
Eq.~\eqref{eq:resource}; interval accounting computes
Eq.~\eqref{eq:exposure} and the matching gain exposure; and nonnegative
shortest-path search computes
Eq.~\eqref{eq:repair}. A final function evaluates explicitly supplied features
and weights after the consequence calculation.

The implementation accepts measured graphs and outcome maps. It does not
pretend to infer a complete causal world from raw observations. Frame identity,
remaining budget, and clock conditions belong in the supplied state model when
they affect admissibility. A missing frame key means a \emph{known} empty
continuation set; a missing measured quantity must not be passed as an empty set.
The same restriction applies to resource entries omitted as known zero.

For finite graphs, stepwise reachability visits at most all supplied edges at
each of $h$ steps, giving an $O((h+1)(|V|+|E|))$ bound including materialisation, before
outcome processing.
Hash-based set comparison is linear in the represented outcome sets in the
usual expected case. Repair search uses Dijkstra's algorithm with nonnegative
scalar edge costs. Augmenting states with budgets or enumerating long courses
can greatly increase the graph: calculable does not mean cheap at every scale.

Finite adjacency iterables are materialised so that repeated visits preserve
the supplied transition relation. The numerical implementation uses finite
arithmetic: out-of-range computed masses, exposures, scores, or explored
repair-path costs raise an error. They are not reported as measured infinite
burdens or as evidence that repair is impossible.

The functions realise the finite support, resource, path, and repair
specialisation. Access-price profiles use repair
search separately for each target outcome; path accounts across uncertain
courses require those courses to be supplied. A stochastic channel requires
channel-capacity machinery beyond the deterministic function. These extensions
change the measurement implementation, without introducing a moral eligibility
condition.

\section{What follows from the construction}\label{sec:results}

The construction earns its main claims under explicit conditions. In particular,
observer independence concerns whether the consequence occurs; it is stronger
than leaving a weight variable out of a formula.

\begin{proposition}[Independence from passive evaluation]\label{prop:observer}
Fix the actual affected transition relations, including any recorded stochastic
laws or control channels, starting states, expenditures,
and continuity conditions. Add, remove, or change an evaluator whose observation
and judgement do not causally alter those relations. The removed and introduced
continuations and incurred expenditures are unchanged. Where the measurement
contract is also held fixed, $\C_\theta(a\mid b)$ is unchanged.
\end{proposition}
\begin{proof}
Occurrence is determined by the affected starting states and admissible
transitions. By hypothesis these inputs and the expenditure records do not
change with evaluation. Hence the sets of admissible paths do not change,
their aligned outcome images do not change, and their set differences do not
change. The resource differences and course profiles likewise retain their
inputs. No report or judgement is an argument of these operations.
\end{proof}

A closed exit remains closed without a witness. That is the causal content of
the proposition. An evaluator who opens the exit, imposes a restriction, or
pays for an option is an additional actor; the causal comparison has changed.
An evaluator who uses a coarser outcome map changes their account's resolution.
Neither case makes recognition the source of the earlier cost.

\begin{proposition}[Independence from decision weights]\label{prop:weights}
For fixed $\theta$, $a$ and $b$, replacing an evaluation rule $V_w$ by $V_{w'}$
does not change $\C_\theta(a\mid b)$.
\end{proposition}
\begin{proof}
The account uses transition, outcome, expenditure and comparison inputs.
Decision weights occur only in the subsequent evaluation map
$V_w:\C_\theta\mapsto\text{decision score or ordering}$. Replacing that map
leaves its argument fixed.
\end{proof}

\begin{proposition}[Representation invariance]\label{prop:representation}
Suppose two groundings are related by bijections preserving affected identity,
starting states, admissible transitions, consequential outcome identity,
resource incidence, units and numerical expenditures, time alignment, outcome
measures, access-price profiles, and any recorded stochastic laws or control
channels. For the repair component, also preserve
admissible repair transitions, their numerical edge costs, and restoration
predicates. Their
consequence accounts agree under those bijections.
\end{proposition}
\begin{proof}
A transition-preserving bijection carries each admissible path to an admissible
path of the same declared length. Its aligned outcome bijection therefore
carries reachable sets to reachable sets and commutes with set difference.
Measure preservation gives equal losses and gains. Resource and time
preservation give equal burdens and profiles. Repair paths correspond with
equal expenditure and restoration conditions, so their infima agree.
\end{proof}

This states what substrate independence buys: the procedure and its results
survive changes of representation that preserve the relevant structure. Two
embodiments with different resource demands or continuations need not have equal
accounts merely because both are called agents. A simulation label cannot by
itself discount an isomorphic consequence.

\begin{proposition}[Termination remains in the account]\label{prop:termination}
If an affected identity has a nonempty baseline reachable set and no continuing
action outcomes, its lost set equals the baseline set. Its loss remains in the
account even if the post-action population omits that identity.
\end{proposition}
\begin{proof}
The comparison population is the union across branches. For that identity,
$\R_i(a)=\varnothing$, so
$\Lost_i=\R_i(b)\setminus\varnothing=\R_i(b)$. The account is indexed by the
comparison population rather than by the surviving census.
\end{proof}

A copy contributes a continuation of the original only where the continuity
conditions support it. The framework thus makes replacement a question about
actual retained relations. It cannot erase a loss through renaming its bearer.

\section{Decisions and relative weightings}\label{sec:decisions}

A derived cost does not force one action in every conflict. Protecting a recipient
can require restricting an attacker; preserving a system can require spending
resources that would otherwise remain available elsewhere. The account exposes
what each choice changes. A decision rule makes its priorities inspectable.

For numerical components in a declared evaluation space, one example is
\begin{equation}\label{eq:decision}
V_w(a\mid b)=\sum_i\alpha_iL_i-\sum_i\beta_iG_i
 +\sum_{i,r}\gamma_{ir}K_{ir}+\sum_i\delta_iD_i-\sum_i\eta_iD_i^+,
\end{equation}
with an action selected by minimising $V_w$ among feasible alternatives.
Weights express relative priorities and any chosen conversions. When an outcome
identity matters independently of its count, replace the corresponding scalar
by an outcome-level feature vector. Lexicographic rules and constraints can
also evaluate the same account. A hard constraint is a decision condition;
it is not evidence that the excluded action had no measurable consequences.

Resources, capability counts and capability-time exposure cannot be added as
physical quantities without conversion. Equation~\eqref{eq:decision} is a
declared evaluation function. A conversion grounded in a common physical
quantity differs from a stipulated priority between unlike quantities, as
Paper 04 explains (Remark~\ref{NA-rem:kappa-kinds}). Both can be used if their
status is stated. Neither silently changes the underlying quantities.

\begin{proposition}[Conditional dominance]\label{prop:dominance}
In a common declared feature space, let losses, burdens and exposure have
positive signs, and gains negative signs. If $z_j(a)\leq z_j(a')$ for every
feature $j$, then $w\cdot z(a)\leq w\cdot z(a')$ for every nonnegative $w$.
The inequality is strict if a strictly smaller feature receives positive weight.
\end{proposition}
\begin{proof}
Each term $w_j(z_j(a')-z_j(a))$ is nonnegative. Their sum has the stated sign,
and is positive under the strictness condition.
\end{proof}

This is agreement within a class of evaluators, not a proof that every possible
agent must use nonnegative weights. For a declared admissible weight set $W$,
robust preference for $a$ over $a'$ means
$\sup_{w\in W}w\cdot(z(a)-z(a'))\leq0$. If the inequality fails, the disagreement
can concern weights even when every evaluator accepts the measured account.

Consent has a place in both layers. It can causally change coordination,
permission, enforcement cost, and actual reachability. It can also receive
weight in a decision. It does not erase a remaining restriction by changing
its name. Similarly, inability to protest does not remove a bearer from the
account. Paper 01 already distinguishes occurrence of harm from its recognition
(Remark~\ref{GE-rem:admissible-sign}); this construction makes the distinction
explicit in the calculation.

\section{Worked cases}\label{sec:cases}

\subsection{Protection with a calculable restriction}\label{sec:protection}

Let frame $A$ have six baseline outcomes $\{1,2,3,4,5,6\}$ and frame $B$ one
baseline outcome $\{x\}$. A restriction leaves $A$ with $\{1,2,3\}$ and gives
$B$ the outcomes $\{x,y,z,u\}$. Under counting measure,
\[
L_A=3,\quad G_A=0,\quad L_B=0,\quad G_B=3.
\]
Suppose the action also imposes one resource unit on $B$, while the baseline
imposes none. Then $K_{B,r}=1$. A declared decision score is
\begin{equation}\label{eq:example}
V_w=3\alpha_A+\gamma_{B,r}-3\beta_B.
\end{equation}
For $(\alpha_A,\gamma_{B,r},\beta_B)=(1,1,2)$ the score is $-2$; for
$(2,1,1)$ it is $4$. The baseline has score zero in both accounts. The two
evaluators choose differently while agreeing on every measured consequence.

The appendix computes these values by set comparison, rather than defining
loss as a preference. For deterministic controllable channels, $A$ loses one
bit of capacity, $\log_2(6/3)=1$, while $B$ gains two bits,
$\log_2(4/1)=2$. Counts and bits are alternative summaries with explicit units;
neither is added to the other as though it were another independent loss.

Equal counts can conceal a different case. Replacing $\{1,2,3\}$ with
$\{4,5,6\}$ leaves capacity and cardinality unchanged while recording three
lost and three gained outcomes. A rule that uses only net count ignores that
substitution. The fuller account makes the omission visible.

\subsection{Short-term options and continuing viability}

An action can open many immediate options by exhausting a budget required for
later continuation. A short horizon registers expansion; a longer horizon
registers contraction. This is a change of comparison question, not an
inconsistency in the set operations. Include remaining resources in the state
and publish the horizon profile when short and long accounts differ.

The same applies to delegation and automation. A recipient can gain access to
an outcome while becoming dependent on a provider whose withdrawal closes it.
The continuation graph contains that dependency. Counting only the currently
available menu omits the provider relation and the cost of maintaining access.

\subsection{Restoration without cancellation}

Suppose the baseline maintains four capabilities through three unit-length
intervals. The action permits only one through the first two intervals and
restores all four in the third. The loss profile is $(3,3,0)$, terminal loss is
zero, and exposure is $D=6$ capability-time units. If the restoration consumes
two resource units, that expenditure is a separate entry.

Conversely, an action can temporarily add one capability to the same baseline
through the first two intervals, then return to the original four. Its gain
profile is $(1,1,0)$ and gain exposure is $D^+=2$ capability-time units, despite
zero terminal gain. The function \texttt{gain\_profile} calculates this matching
course account.

The account establishes temporary restriction, exposure and expenditure. Whether
these quantities correspond to pain, distress, or another experienced effect
requires a grounding for that effect. A conscious preference to remember the
restriction is unnecessary for recording the structural loss. Byte-identical
restoration does not make the two courses identical.

\subsection{Termination and replacement}

A frame with baseline continuations $\{r,s,t\}$ is terminated in the action
branch. Its continuing outcome set is empty and its loss is three. A newly
created frame with three outcomes does not cancel the loss by cardinality.
If a repair or transfer preserves the original's identified relations, the
continuity contract can establish which outcomes remain available. The account
then changes because retained structure has changed, rather than because a
new entity was announced as equivalent.

The same rule applies to biological and computational cases. Their expenditures
and continuity relations can differ substantially. No category supplies a
blanket exemption or a universal equality of numerical costs.

\subsection{Disposable encounters and persistent companions}\label{sec:games}

Games offer a particularly clear comparison because their continuing relations
can be inspected in the program. A player can treat the loss of one companion
as more consequential than the removal of a thousand encounter enemies. The
framework distinguishes two mechanisms: the underlying structures can lose
different futures, and a player can weight their recorded losses differently.
The first mechanism exists before the second.

Consider a finite game with an encounter character $e$. In both branches the
encounter ends and the character has no persistent state, later action, social
relation, or continuing world role. Let the comparison concern the state after
that boundary. The baseline resolves the encounter without the depicted killing;
the action resolves it with the depicted killing. If no continuing outcome of
$e$ is instantiated in either branch, then
\begin{equation}\label{eq:npc}
\R_e(b)=\R_e(a)=\varnothing,\qquad L_e=G_e=0.
\end{equation}
The zero concerns actual absent continuing futures. It is not a low weight
assigned to a rich persisting agent. If the wider world, player state, and
resource expenditure also agree, those account entries agree as well.

Now let a companion $c$ have three distinct reachable continuing outcomes: a
later joint action, an ongoing exchange, and a persistent world response.
Specify $\R_c(b)=\{r_1,r_2,r_3\}$ and $\R_c(a)=\varnothing$ under an outcome
map that distinguishes these possible continuations. Then
$\Lost_c=\{r_1,r_2,r_3\}$ and $L_c=3$ under counting measure. This toy count
concerns distinct outcomes, rather than counting three concurrent features of
one outcome as if they were separate channel outputs.
Effects on a player and other characters belong to their corresponding accounts.
They need not first judge the death meaningful for those continuations to be
removed. Player valuation explains an additional part of the resulting response.

Equation~\eqref{eq:npc} requires its stated boundary. An encounter character can
have intermediate options, a restriction interval, or effects on another agent.
If killing it early removes those options, its path account can be positive even
when its terminal continuing account is zero. If it learns, persists, or sustains
relations outside the encounter, its future set is different. ``Designed to be
killed'' is therefore evidence about a role, not a mathematical proof of zero
cost. The code and actual continuation relations settle the case.

Likewise, a depiction can carry substantial consequences in a player's
relationships, retained game history, or other instantiated structures even if
the target has no persistent organisation. A represented life story can become
part of such a structure. The account follows the relations that exist rather
than granting or removing consequences through the adjective ``fictional.''

A thousand characters satisfying Eq.~\eqref{eq:npc} do not create a thousand
lost continuing lives merely because the depicted acts share the label killing.
In a common game-world account with identical continuing outcomes, the loss
remains zero. Positive intermediate or resource entries, where present, remain
positive; aggregation requires their common units and incidence. This explains
how the informational cost of a disposable encounter can be nearly trivial while a
companion's loss is substantial.

The 1990s video-game controversy supplies historical context. The US Senate's
activity report documents the 1993 and 1994 hearings on violent game content
\citep{senate1995activities}. A content label identifies what is depicted. The
present theory asks which actual futures, relations, and resources the depicted
act removes. The two classifications need not coincide. Research on game-related
guilt and moral disengagement addresses players' responses to such framing
\citep{hartmannvorderer2010disengagement}; that response is a measurement or
valuation phenomenon, not proof of the target's underlying continuation account.

\section{Moral injury as a whole-system consequence}\label{sec:moral-injury}

Moral injury shows why an account confined to an action's external recipient
is incomplete. The apparent extra ethical cost is an omitted internal or
relational cost of the system. An actor can alter their own continuing organisation by doing
something, being unable to prevent it, or receiving information about it.
Witnesses and institutional relations can also bear persistent consequences.
The resulting costs need not disappear when the action receives a favourable
ethical judgement.

Litz and colleagues describe lasting consequences following events that violate
deeply held moral expectations, including perpetration, failure to prevent, and
witnessing \citep{litz2009moralinjury}. Shay emphasises high-stakes betrayal by
legitimate authority and impaired trust \citep{shay2014moralinjury}. These accounts
supply the phenomenon to explain. The present framework proposes an
informational interpretation of its changes, rather than substituting a
capability count for a clinical construct.

\subsection{Self-relations and trust are enabling structure}

A person's expectations about themselves and others can enable continuations:
accepting help, cooperating, making commitments, exercising a role, or relying
on an institution. These possibilities depend on actual embodied, cognitive,
and social organisation. If an event changes that organisation, it can close
continuations or increase the effort needed to sustain them. The account uses
the same loss, access-price, expenditure, and path functions already defined.
No extra moral-injury eligibility rule is needed to record the change.

Unchanged support sets do not establish unchanged reliability or burden. Where
trust changes the probability of successful cooperation while leaving it possible,
retain the grounded transition probabilities or control channel alongside the
support account. Where an experienced or physiological burden is instantiated,
its own grounded dimension is required. The finite set functions alone do not
certify the absence of those effects.

This is the proposed mechanism at the structural level. The event changes
retained information and enabling relations; those changes alter subsequent
admissibility and costs. A self-model or expectation matters because of what it
causally enables or restricts. Mere disagreement between two reported beliefs
does not establish that a capability was lost. Conversely, a damaged relation
can matter even when its bearer cannot articulate an ethical explanation.

Internal moral commitments can enter this causal process without becoming
externally imposed accounting axioms. They are part of the organisation of the
particular system being studied. A commitment violated by an event can change
that system's future behaviour, relations, and expenditure. Another system with
different commitments can undergo different changes. The functions remain the
same while their grounded inputs differ.

\subsection{A chosen action can retain an internal bill}

Extend the protective example of Section~\ref{sec:protection} with an actor $P$.
Let its baseline continuing outcomes be
$\R_P(b)=\{u,v,w\}$ and its post-action outcomes be $\R_P(a)=\{u\}$, where the
outcome map distinguishes grounded continuations involving cooperation and
participation. If the action also imposes two additional resource units on $P$,
then $L_P=2$ and $K_{P,r}=2$. These are illustrative consequences of the specified
model, not numerical findings about a person.

The original protective account remains: $L_A=3$, $G_B=3$ and $K_{B,r}=1$.
The expanded decision score includes
\[
V_w=3\alpha_A+\gamma_{B,r}-3\beta_B
       +2\alpha_P+2\gamma_{P,r}.
\]
For weights that place sufficient priority on protection, this action can
still be selected. Its selection does not cancel the actor's loss. The full
account explains how an action regarded as necessary or justified can carry a
persistent cost for the system that performed it. Establishing moral injury in
a particular person requires evidence about the event, the affected organisation,
and its consequences; the toy account supplies no diagnosis.

The same reasoning covers a witness whose continuing trust is damaged, or a
person unable to prevent an event. Attribution asks what actually changed and
through which relations. It does not require the affected person to have caused
the original external injury. An institution's practical ability to coordinate
can likewise change when trust relations are damaged. Individual and collective
entries must retain their overlap rather than be indiscriminately added.

\subsection{The trolley switch: external consequences and operator state}

The trolley problem makes this separation explicit. Thomson's treatment already
compares the trolley driver and a bystander at the switch, including their
responsibilities and discomfort \citep{thomson1985trolley}. The present account
uses the established dilemma to calculate separately attributed consequences.

Fix $a$ as operating the switch and $b$ as leaving it unchanged, with the usual
stipulated external outcomes: $a$ saves five people and kills one who would
survive under $b$. For each saved person $k$, retain the grounded gain
$G_k(a\mid b)$; for the person killed, retain $L_6(a\mid b)$. The headcount
specifies neither the lost continuations nor a universal exchange rate between
these persons. Their identified futures, external resource changes, witnesses,
and institutional effects belong in the external account under its declared
boundary.

Let $x_P$ describe the operator's initial organisation, including commitments,
self-model, experience, knowledge, and enabling relations. Write the operator's
attributed component as
\begin{equation}\label{eq:trolley-operator}
\C_{\theta,P}(a\mid b;x_P).
\end{equation}
This is calculated with the same functions, using the operator's branch-specific
continuations, expenditures, path profiles, and repair conditions. It is not a
casualty count converted into guilt. Operating the switch can change one
operator's continuing self-relations or impose substantial expenditure, while
another operator undergoes different changes under the same external outcome.
Declining to act can likewise impose an internal burden; the baseline is not
presumed free of consequences.

For a finite illustration, fix the same external branches and give two possible
operator states the baseline outcome set $\{u,v,w\}$. In state $x_P$, switching
leaves $\{u\}$ and adds two resource units, yielding $L_P=2$ and $K_{P,r}=2$.
In state $x'_P$, switching leaves $\{u,v,w\}$ with unchanged expenditure,
yielding $L_P=0$ and $K_{P,r}=0$ in these dimensions. Other dimensions remain
to be checked. These are specified toy inputs, not measured psychological
values. They demonstrate that operator cost is calculable and state-dependent
while the fixed external losses and gains remain the same.

The separation is one of attribution, not causal isolation: what the operator
knows or believes about the external event can change their internal transition.
Nevertheless, the external impact alone does not determine the operator's cost.
Their feeling of justification cannot erase another person's loss, and their
internal burden cannot be inferred solely from the number saved or killed.
A decision can therefore favour switching while retaining a substantial
operator cost. Relative priorities decide how these separate entries enter the
choice; no single preferred answer to the dilemma follows from accounting alone.

\subsection{Psychopathic traits and differences in ethical aperture}

The same separation permits an account of differences associated with
psychopathic traits. Aperture retains its inherited meaning: the input gate
and its represented dimensions, resolution, and extent (Paper 01,
Definition~\ref{GE-def:aperture}; Paper 02, Remark~\ref{SM-rem:aperture}).
An \emph{ethical aperture} is the application of that structure to
ethically relevant distinctions and their fit to acceptable future shapes
(Paper 02, Remark~\ref{SM-rem:aperture-shape}). It is a shape, rather than one
quantity of ethical sensitivity. Different frames can represent different
bearers, kinds of consequence, degrees of distinction, and horizons. An ethical
cost dimension absent from a frame is undefined there, rather than present
with value or weight zero. This is absence to that frame, not absence of the
affected relation from the world.

The shape of what remains acceptable adds the flexibility described by this
application. Let $\widehat{\mathcal S}_j(x)$ be actor $j$'s represented feasible
continuations from state $x$. Where the actor holds an operative ethical
acceptability predicate $\Gamma_{E,j}$, define
\begin{equation}\label{eq:ethical-region}
\mathcal E_j(x)=\{y\in\widehat{\mathcal S}_j(x):
                    \Gamma_{E,j}(y\mid x)\}.
\end{equation}
This is a conditional filter on the represented manifold. It is distinct from
the input aperture that determines which distinctions can enter that manifold,
and from the subsequent weights ordering its members. It is not a moral
eligibility gate in the consequence account.

Holding the represented manifold fixed, more permissive ethical conditions can
admit more future shapes; readily revised or context-dependent boundaries can
make the operative ethics flexible. An actor may instead retain only a few
such conditions, or no effective ethical restriction over the compared domain.
In that last case the ethical filter excludes nothing there: effectively
$\mathcal E_j(x)=\widehat{\mathcal S}_j(x)$. This does not require an ethical
predicate to be explicitly represented when none is operative, nor make physical,
resource, or other nonethical constraints disappear. Comparative ethical
priorities can also remain weak or strong within the admitted region. Few
constraints, low weights, and absent dimensions are different possible structures.

Under the inherited fit/adaptive-resolution hypothesis (Paper 02,
Hypothesis~\ref{SM-hyp:adaptive-resolution}), an already acceptable continuation
can attract little further ethical scrutiny: it passes with room at the represented grain.
A continuation near a held boundary can instead attract refinement or pressure
to reshape the acceptable region. What lies outside a represented dimension or
horizon never reaches that fit comparison. These structural possibilities and predictions alone do not establish a
globally wider aperture or a generally more flexible person; comparisons must
state the dimensions, horizons, and held conditions. The application uses
specified trait profiles under matched conditions; psychopathy or sociopathy
labels alone do not supply these inputs.

On this hypothesis, the patterns grouped under psychopathy can be described as
a characteristic family of ethical-aperture shapes and acceptable-region
boundaries. Membership in that family is an empirical question, assessed through
represented distinctions and held constraints across situations. The hypothesis
does not identify every such profile with one identical shape.

The proposed definition is intrinsic: its defining properties are the person's
operative distinctions, acceptability boundaries, and their dynamics within
that person's actual relational organisation. They exist whether or not an
observer identifies them or applies a diagnostic label. Intrinsic does not
mean independent of context: learning, social coupling, and current conditions
can constitute or change that organisation. A report or classification describes
the structure under a measurement contract; it does not confer the structure.

Practical recovery is an inverse problem. Similar observed choices can arise
from different apertures, acceptable regions, weights, beliefs, or constraints,
and a self-description need not identify the operative profile. Matched
comparisons may narrow the compatible family without recovering one unique
shape, as in the conditional reconstruction of Eq.~\eqref{eq:revealed-self}.
The distinction between a defined structure and uncertain access to it is
therefore retained: uncertainty about the profile belongs in the measurement,
without making the defining properties dependent on an observer's judgement.

A general account of psychological differences would be a separate programme.
Many clinical categories are operationalised through symptoms, distress, or
impairment; NIMH's RDoC framework instead investigates underlying functional
dimensions and mechanisms that can account for symptoms \citep{nimhrdoc}.
This is close prior work for the proposed direction. In the present proposal,
felt difficulty is a modelled assessment of an operative state or transition.
The intrinsic definitional object is the operative organisation, independently
of whether its model assesses that organisation as difficult or problematic.
Such assessments can provide evidence and causally change subsequent conduct;
they need not identify the underlying profile. The relevant distinction is
between diagnostic recognition criteria and a defining account of operative
organisation. The present paper
uses ethical-aperture differences to illustrate consequence accounting. It does
not develop a general psychological classification or establish that its
structural profile replaces a clinical diagnosis.

The actual consequence account remains distinct from this aperture. A victim's
lost continuation does not become smaller because an actor represents it poorly
or gives it little weight. Conversely, reduced self-censure, aversion, or
identity conflict can be a real difference in the actor's own consequences,
where grounded in their organisation. This is a difference in attributed cost,
in addition to any difference in its evaluation. It connects directly to the
state-dependent operator account in Eq.~\eqref{eq:trolley-operator}.

For a conditional representation, let $o_j$ be the information available to
actor $j$ and $Q_j$ its processing into admitted features. The represented
features are $\widehat z_j=Q_j(o_j)$, and a stated linear decision model applies
$w_j\cdot\widehat z_j$. Different input apertures or $Q_j$ can produce different admitted accounts
from the same available information. Different ethical-region predicates can
admit different continuations within a common represented manifold; different
$w_j$ can produce different choices within the same admitted region. Different actual actor states can also
produce different internal burdens. These are separately investigable routes.
A difference in behaviour alone does not identify which route changed, nor
demonstrate that every difference is a fixed weighting. Attention, learning,
impulsivity, and policy stability can require additional state or decision terms.

Relevant research distinguishes abstract judgement from hypothetical choice.
In a study of 102 students, Tassy and colleagues found that elevated psychopathic
traits predicted more sacrificial, outcome-favouring answers to hypothetical
choice questions, without the corresponding association for acceptability
judgements \citep{tassy2013choice}. The result concerns self-report traits and
imagined choices, not enacted harm. Hauser and colleagues studied 60 male
offenders and found distinct associations by psychopathy facet and
emotion-processing condition \citep{hauser2021subcomponents}. Conditions also
used different dilemmas and were not randomised, limiting causal attribution
to emotion processing. These findings support separating
explicit judgement, processing, and choice, while qualifying a uniform profile.
They do not recover the framework's aperture geometry, acceptability predicates,
or weights, and do not establish absent ethics, zero empathy, or zero internal
cost for every person with elevated traits.

Under this interpretation, differences in ethical application are differences
within the informational and decision machinery already specified. The account
can register the actor's smaller or different internal burden alongside unchanged
external loss. It does not require an extra category of existence or an eligibility
exception for those affected. This is an explanatory application of the framework,
with its empirical mechanisms still to be identified under matched conditions.

\subsection{Recognition can cause further cost without creating the original cost}

Observer independence does not require every consequence to precede awareness.
An external recipient's lost future can exist before anyone knows about it.
Later learning about the event can transform an actor or witness, creating
additional costs in that system. The information received is then part of the
causal transition, rather than a passive observer conferring ethical value.
Proposition~\ref{prop:observer} holds for the fixed causal case; changing what
information reaches a system changes that case.

A false belief can therefore have real consequences for its bearer even when
the supposed external event did not occur. Record the represented event's
causal effects in the bearer and distinguish them from the uninstantiated
external injury. Conversely, an external loss is not removed because no actor
feels guilt. Guilt is not an accounting instrument that certifies or nullifies
all other costs.

Repair has the same separation. Restoring an external recipient's capability
does not automatically restore the actor's self-relations, a witness's trust,
or institutional coordination. Each requires its own grounded restoration
condition and feasible repair paths. Nor does later repair cancel the intervening
path account. Whole-system accounting makes these different bills visible in
one typed family while keeping their incidence and units distinct.

\section{Operative valuation and a self-model that rejects it}\label{sec:numeric}

Consider a self whose decisions already implement a numerical valuation or its
functional equivalent, while its self-model rejects the explicit account of that
valuation as uncomfortable. It can hold that a human life cannot have a value
assigned while its operative choices still make trade-offs involving lives.
Writing an actual number exposes the valuation to the self-model. The distinction
is between what the self carries and what its model of itself admits.

Paper 01b supplies this distinction: the self-model is a proper, incomplete
part of the self, and its realisation account distinguishes changed modelling
from newly occurring facts (Section~\ref{MI-sec:selfmodel},
Remark~\ref{MI-rem:realisation}). An operative valuation need not be stored as
a literal numeral. It can be realised through relative weightings, thresholds,
or an action ordering. The present hypothesis concerns a valuation already
operative in those relations. Explicit acknowledgement need not create it.

The closest prior work includes taboo trade-offs, sacred values, and motivated
reasoning. Fiske and Tetlock examine why some explicit exchanges are offensive
in themselves \citep{fisketetlock1997taboo}. Tetlock and colleagues studied
responses to monetising protected values \citep{tetlock2000unthinkable}, and later
work discusses shadow prices under scarcity \citep{tetlock2017sacred}. Kunda's
account describes motivation shaping access to and evaluation of beliefs
\citep{kunda1990motivated}. These supply relevant mechanisms and observations.
The specific explanation proposed here is that discomfort can protect a
self-model from admitting an already operative valuation. Discomfort alone is
not proof that this explanation applies in a particular case.

\subsection{Models of an agent have no automatic accuracy hierarchy}

Can another person know an agent better than the agent knows themselves? For
specified aspects of the agent, the answer can be yes. Paper 01b's self-model
account distinguishes access mode from accuracy: the self-model is a proper,
incomplete part of the self, while an external model has its own access,
resolution, and modelling investment (Proposition~\ref{MI-prop:selfmodel} and
its subsequent discussion). Neither position alone establishes greater accuracy.

Impaired self-awareness supplies a concrete case. Cocchini and colleagues
documented unawareness of motor impairment after left-hemisphere brain damage,
using neurological assessment, self-assessment, and caregiver ratings
\citep{cocchini2009anosognosia}. Observed impairment can conflict with the
person's account of their own ability. The study also illustrates that the
assessment method affects what can be recovered: language-dependent interviews
and a less language-dependent instrument identified different levels of
unawareness. This concerns particular impairments and assessments, without
making disability itself a general loss of self-knowledge.

Self-modelling and modelling other agents are closely related capabilities.
Both construct selective representations and can use shared skills in
distinction, inference, prediction, and calibration. Their embedding differs:
a self-model is an operative part of the organisation it models, so revising it
can change its own target's subsequent dispositions and continuations. An
external model has its own access and affects the target through whatever
coupling is instantiated. Skill in one can support skill in the other, while
their different access and embedding prevent accuracy in one from guaranteeing
accuracy in the other.

Domain expertise provides a further illustration. A master may recognise a
junior's likely move, mistake, or response from a familiar configuration before
the junior can predict or articulate it. The master's learned structure can
resolve domain distinctions and continuation patterns that the junior's
self-model does not yet hold. The comparison concerns predictive accuracy
within that domain under matched information and conditions, assessed against
the junior's subsequent actions. It need not imply broader access to the
junior's experience or greater accuracy about the whole person. This is an
illustration of how modelling skill can affect comparative accuracy, rather
than a new empirical result of this paper.

Every model of an agent is therefore examined through its own aperture and
calibration, including the agent's self-model, another person's model, and an
institution's model. Aperture specifies represented dimensions, resolution,
and horizon; accuracy concerns agreement with grounded observations or
predictions under declared comparison conditions. A model with broader coverage
can be poorly calibrated, while a narrower one can be accurate within its
scope. Models can also be accurate on different aspects without one being
uniformly better. Accuracy comparisons require aligned target properties,
conditions, and a stated criterion.

For ethical inquiry, self-report is evidence about a model of the operative
self, rather than a definition of the self's actual capabilities or dispositions.
External models can improve that account where their comparative accuracy is
grounded. They can also omit or misrepresent it. This preserves the intrinsic
definitional object while allowing variable access to it, and explains why
recovering an operative ethical profile can require comparisons among several
models instead of accepting one vantage as definitive.

\subsection{Recovering the operative self from choices}

A fundamental epistemic problem of ethical inquiry is the relation between the
self and its self-model. The working hypothesis is that much ethical inquiry
attempts to reconstruct the self's operative dispositions from revealed choices,
and to calibrate its self-description against that evidence. Here ``recovering
the self'' means recovering aspects of its consequential organisation; finite
observations need not identify the entire self or a unique internal mechanism.
The self-model is itself active within that organisation: its commitments can
change later choices, rather than merely report an otherwise fixed actor.

Revealed-preference theory supplies an established method for inferring
comparisons from choices. Samuelson's original construction compares a selected
bundle with an alternative that was affordable under the same conditions
\citep{samuelson1938behaviour}. The present application compares inferred operative
dispositions with the actor's admitted self-description. It does not require a
literal numerical representation inside the actor.

For a conditional linear example, let observation $k$ record the selected action
$a_k$, the alternatives $A_k$ actually available and known to the actor, and
aligned feature vectors $z_k(a)$ under the actor's outcome beliefs. Let $W_0$
be a declared family of candidate weights, with units and any normalisation
specified. For informative weight reconstruction, exclude the zero vector,
which otherwise makes every action tie, and declare a scale convention if a
unique numerical magnitude is sought: positive rescaling preserves these
choice inequalities. Assume that one stable weight vector minimises the score
in all recorded comparisons, allowing ties. The compatible operative models are
\begin{equation}\label{eq:revealed-self}
W_{\mathrm{op}}=
\{w\in W_0:\ w\cdot z_k(a_k)\leq w\cdot z_k(a)
\quad\text{for every }k\text{ and }a\in A_k\}.
\end{equation}
On a finite candidate family this is an implementable filter using
\texttt{evaluate}. Equation~\eqref{eq:shadow} below is its one-conversion case.
Additional comparisons can shrink the compatible family while the model and
conditions remain fixed. An empty family rejects the stated candidate family
or choice assumptions; several surviving models leave the reconstruction
underdetermined. A stochastic or changing policy requires a correspondingly
different inference model.

Feasibility, knowledge, coercion, outcome beliefs, and context matter to this
inference. An unavailable alternative supplies no rejected-choice evidence.
Features inferred from an actor's beliefs need not equal the measured actual
consequences: both belong in the account with their distinction retained.
Revealed preferences help reconstruct how the actor responds; the affected
systems' losses do not depend on that reconstruction. Nor does discovering an
operative disposition establish that it ought to be retained. Inquiry can
calibrate the self-model, revise the operative disposition, or change the
conditions of action. These are distinct interventions with calculable effects.

\subsection{Energy-state management and deformation of decision fields}

Within the proposed ontology, the foregoing processes concern states and
transitions in an $n$-dimensional relational manifold. Relevant dimensions can
include resources, enabling relations, affective burdens, self-model commitments,
social sanctions, and the continuations available to affected frames. Paper 02
supplies the operational manifold, local transition costs, and deformability of
the fields over it (Definitions~\ref{SM-def:seldon}, \ref{SM-def:contours}, and
\ref{SM-def:deformability}). The self and its model are part of the same evolving
organisation, so changing the model is itself a transition in that manifold.
It can also deform the decision field governing subsequent transitions.

For a conditional scalar decision model, write
\begin{equation}\label{eq:model-field}
\Phi_M(a\mid x)=w_M\cdot\widehat z_M(a\mid x),
\end{equation}
where $x$ is the current joint state, $\widehat z_M$ records the action features
admitted by the model, and $w_M$ supplies declared conversions into the score's
unit. Under a stated minimising policy, reducing an action's score relative to
its alternatives can change which transition is selected. A change in represented
loss, self-censure, anticipated social sanctions, or expected benefit can all
change this ordering. Energy and resource quantities retain their grounded units;
a single scalar decision field exists only after its conversions are specified.
Capacity loss remains a separately typed consequence even where energy is spent
to produce or prevent it.

This supplies a common geometry for apparently different ethical moves.
Admitting an uncomfortable valuation can require a costly self-model transition.
Keeping the valuation unadmitted can preserve an existing organisation while
constraining later modelling. A forced comparison can make that organisation
incompatible with the requested commitment. Reclassifying an affected group
can lower the represented cost of harming its members. Each is a possible way
of changing states or transition prices in the manifold. The claim that organisms
implement these moves through compressed regulation is the paper's explanatory
hypothesis, rather than a recovered neural energy model.

The two continuations ``act under the current model'' and ``change the model,
then act'' must therefore both be accounted for. The second includes the price
of changing the field: propaganda, social coordination, and internal revision
need not be free. Once a category rule is operative, it may apply across many
individuals without a separate confrontation with each person's relations.
Whether that makes the complete course cheaper for a particular actor requires
the declared account. Making a subsequent action cheap in a decision field does
not imply that its affected frames lose fewer futures.

\subsection{Play as an elicitation of otherwise unspoken comparisons}

Games such as ``fuck, marry, kill'' and ``would you rather'' provide familiar
illustrations of deliberately eliciting comparisons. The first assigns three
people to sharply different imagined relations; the second requests a choice
between specified alternatives. The first compares assignments of relations,
which need not give one total ranking of the people across those relations. A participant can discover an ordering they
had never made explicit, or encounter a comparison that their self-description
resists acknowledging. The proposed mechanism is a forced local collapse of
self-model ambiguity: answering commits the participant to a comparison that
could previously remain open. In this series' terms, collapse resolves a
particular distinction under a constraint; it need not destroy the self or
settle every aspect of its model. The game makes continued noncommitment
incompatible with supplying the requested answer.

Let $D_M$ be the descriptions the self-model leaves open before the question,
and let $D(q)$ be those compatible with an answer $q$ under its literal reading.
A self-model that admits the answer as evidence restricts its open descriptions
to $D_M\cap D(q)$. A proper nonempty subset records a local resolution of
ambiguity. An empty intersection exposes a conflict requiring revision of the
prior description or of the admitted meaning of the answer. The forced answer
creates a commitment to confront; psychological admission of that commitment
remains a further event. A participant can refuse, joke, or reject the answer's
application to themselves. Those responses locate where the attempted collapse
was resisted, rather than proving that the implied self-description was admitted.

An ordering is already a functional equivalent of valuation for those
alternatives, although it need not identify a cardinal number or a stable ranking
in other contexts. The mechanism hypothesis is that playful forced choice can
make an operative comparison explicit and pressure the self-model to resolve
its account of that comparison.

These are hypothetical elicited preferences. An answer is an actual choice of
what to say in the game, with possible social and self-modelling consequences;
it is not an observed performance of the imagined actions. Humour, audience,
provocation, refusal, and the imposed options can shape the answer. The examples
therefore illustrate a method of confronting the self-model with comparisons,
without independently validating the proposed psychological mechanism. A study
would distinguish answers, admission as self-description, and consequential
choices under matched conditions. The game itself can also change the actor,
so later answers need not simply uncover a valuation unchanged by elicitation.

\subsection{Refusing a possible action versus judging it unlikely}

Statements of surprise in crime programmes provide another illustration:
``I cannot imagine someone like that doing this.'' The proposed interpretation
is that some such responses express refusal to admit a reachable action into a
model of oneself or another person, rather than a prior estimate that the action
had low probability. The observed statement alone does not identify that
mechanism. Surprise can also reflect a reasonable low-probability estimate,
lack of relevant knowledge, or the social expression of condemnation.

The formal distinction concerns the model's represented possibilities. Let
$\widehat R_M$ be the action continuations admitted by a self-model or a model
of another actor, and let $R_\theta$ be the independently grounded reachable
continuations under the declared conditions. A particular harmful continuation
$h\in R_\theta\setminus\widehat R_M$ is a model omission. If $h\in\widehat R_M$
and the model assigns it a small probability, it is instead a represented but
unlikely event. If no probability has been assigned, omission is not a measured
probability of zero. The difference is between which alternatives the model
admits and how it distributes expectation among admitted alternatives.
Establishing that an omission is motivated requires evidence beyond establishing
the omission itself.

A model of ``people like us'' can support a person's own self-description.
Evidence that an apparently ordinary person performed a harmful action can
therefore require revision both of the model of that person and of the observer's
account of their own possible conduct. The protective-self-model hypothesis
developed here predicts that some resistance occurs at this admission step;
compression alone does not require resistance to updating. Recognition makes a previously
excluded continuation available to modelling; it need not make that continuation
more reachable in the original event. Recognition can nevertheless change the
observer's later actions and relationships, whose consequences belong in the
whole-system account.

Bandura's account of moral disengagement is relevant prior work: it describes
self-regulation and mechanisms for preserving a moral self-description while
performing harmful actions, including justification, displacement of
responsibility, and minimisation of harm \citep{bandura2002disengagement}.
This supports examining interacting personal and social conditions rather than
assuming a moral identity fixes conduct across circumstances. It does not
establish that every individual can perform every crime, or diagnose why a
particular viewer expresses surprise. The claim needed here is narrower:
self-description alone cannot exclude a continuation whose feasibility has
been independently grounded. Reachability is conditional on actual capability,
access, resources, and circumstances. Possibility also does not establish a high
probability, willingness, or a uniform propensity across people.

\subsection{Dehumanisation and exclusion from the represented account}

Two changes of the decision field should be distinguished. An actor can retain
an entity's represented losses and regard violence against it as deserved or
justified, giving the action sufficient favourable significance in the evaluation.
Alternatively, the actor can lower the represented standing of the entity so
that harming it contributes little loss to the decision field. The proposed
interpretation of the genocidal reclassification move concerns this second
route: reducing the represented cost of harming members of a category, without
requiring a separate judgement that each deserves punishment. The routes can
coexist in the same actor or institution.

An actor who regards killing a person as costly can reclassify a group
as non-persons, pests, or threats, and thereby remove or attenuate the losses
admitted by their own decision model. Group classification can apply that
exclusion across many individuals without confronting each person's continuing
relations. The proposed mechanism concerns a reduction in represented loss and
in restraints on action; it does not reduce the affected person's informational
loss merely by changing a label.

For an illustrative binary exclusion rule, let $e_{i,M}\in\{0,1\}$ specify whether
an actor's model admits an affected frame's loss entry. For a fixed finite grounded
loss $L_i$, its represented entry would be
\begin{equation}\label{eq:category-exclusion}
\widehat L_{i,M}=e_{i,M}L_i.
\end{equation}
Changing $e_{i,M}$ from one to zero erases that entry from the actor's represented
account while leaving $L_i$ fixed. This describes a gate in the actor's model;
it is not an eligibility rule of the present framework. Partial attenuation,
changed weights, and reinterpreted outcome beliefs provide other ways of
changing an operative evaluation. The person whose future is destroyed remains
an affected frame in the whole-system account regardless of the perpetrator's
category assignment.

Kelman's account of sanctioned massacres is particularly close to this
interpretation: it examines authorisation, routinisation, and dehumanisation
as processes weakening restraints on violence against groups that have not
themselves engaged in hostile actions \citep{kelman1973violence}. Landry and
colleagues' analysis of Nazi propaganda found decreasing attribution of mental
experience to Jews before the Holocaust, consistent with disengagement of moral
concern \citep{landry2022propaganda}. The latter measures propaganda language,
not the distribution of individual perpetrators' decision mechanisms; its
observational findings and missing archival data do not establish causal
sufficiency or a majority prevalence. The predominance of the proposed route
across perpetrators remains an empirical question.

Bandura identifies dehumanisation and victim blaming among mechanisms for
selective moral disengagement \citep{bandura2002disengagement}. The proposed
informational reading is that reclassification can permit a moral self-model to
remain compatible with harmful action by excluding consequences or changing
their interpretation. Dehumanisation is not a necessary explanation for every
atrocity. Rai and colleagues distinguish instrumental violence from violence
perceived by its perpetrators as morally justified; their experiments found
increased willingness for instrumental harm under dehumanising descriptions,
without the same increase for morally motivated harm
\citep{rai2017dehumanization}. These studies concern attitudes and vignette
responses, not a complete causal demonstration of genocide. A perpetrator can
recognise a victim's humanity and still construe harming that victim as deserved.
Both exclusion and justification can be analysed as changes to the model or its
evaluation, with their actual consequences retained.

The reclassification process also has consequences of its own. It can reduce
real self-censure or social sanctions for perpetrators, while propaganda,
isolation, institutional exclusion, and organised violence change victims' and
witnesses' actual enabling relations and futures. These effects must be counted
where grounded. The framework therefore permits an actor's burden to fall while
another frame's loss rises. A saving in one attributed dimension does not erase
a loss in another. Relabelling with victims' capabilities and continuations
held fixed does not change their grounded loss. Effects of that relabelling on
actors, institutions, or other affected frames can still change the whole-system
account, even when the victims' measured loss remains fixed.

\subsection{Research support for context-sensitive ethical application}

The proposed mechanism has relevant experimental and historical support.
These findings support context-sensitive choices and changes in admitted
consequences; their interpretation through the present account remains a
hypothesis. They do not independently establish a population distribution of
intrinsic moral character or identify the brain's implementation.

Dana and colleagues experimentally reduced transparency between choices and
recipients' outcomes in a binary dictator game. Generosity fell when actors
could leave consequences uncertain, even though a fair outcome remained
available \citep{dana2007wiggle}. This is relevant to preserving a favourable
self-description while relaxing practical constraints on selfish action.

Vu and colleagues' meta-analysis covered 6,531 participants and 56 treatment
effects. It reported approximately 40\% avoiding easily obtainable information
about consequences for others, and a 15.6-percentage-point reduction in altruistic
behaviour when the opportunity for ignorance was available
\citep{vu2023ignorance}. The result supports information avoidance as a mechanism
for changing conduct. Avoiding new outcome information does not, by itself,
demonstrate refusal to acknowledge an already operative valuation; those are
distinct evidential claims.

In genocide research, Luft's analysis of interviews with 31 Hutu documents
behavioural crossing between killing, desistance, and saving during the 1994
Rwandan genocide \citep{luft2015dynamic}. This supports a dynamic account of
conduct rather than fixed actor categories. The qualitative sample does not
estimate a universal proportion of ambivalent people.

Luft's later analysis examines selected first-time killers whose initial
cognitive and emotional difficulty diminished as killing became routine and
victims became socially meaningless \citep{luft2023career}. The study triangulates
retrospective perpetrator accounts with reports and survivor/rescuer testimony.
It closely matches lowered represented relevance, while qualifying causal order:
in the examined trajectories, dehumanisation followed participation and
habituation. The change need not involve conscious reclassification or explicit
self-talk. This supports possible feedback from action into subsequent decision
fields, without proving that every perpetrator follows that trajectory.

The lack of explicit acknowledgement matters to the self-model hypothesis.
Where an actor's admitted standards condemn a change in conduct, consciously
admitting that change as their own disposition would require revising a
favourable self-description: ``I am becoming someone who does this.'' Implicit
habituation can instead change operative responses and represented relevance
while that description remains uncorrected. In the proposed account, this can
avoid the additional burden of confronting a conflict between conduct and
admitted identity. The actual self changes; its model need not admit the change
as becoming a worse person. Bandura's account supplies closely related prior
work on harmful conduct alongside preserved moral self-description
\citep{bandura2002disengagement}. Conscious reinterpretation can also preserve
that description by presenting the conduct as justified. Thus consciousness
alone does not entail self-condemnation. The specific prediction concerns
acknowledgement of deterioration under the actor's own retained standards.
Lack of explicit acknowledgement of that deterioration need not mean lack of
awareness of the acts or their consequences. Whether avoiding this confrontation
helps sustain a particular trajectory requires evidence beyond habituation
alone.

Asch's conformity experiments supply adjacent evidence that social conditions
change responses while leaving substantial individual resistance
\citep{asch1956independence}. Their task concerns perceptual judgements, not
ethical choice. Asch also reports a non-normal distribution of observed error
counts. Those counts neither establish nor refute a normal distribution of
latent ethical margins; the quantities and tasks differ.

Together these findings give the framework empirically supported phenomena to
explain: information-conditioned prosocial action, changing social relevance,
behavioural switching, and pressure alongside resistance. They support
susceptibility to changes in a decision landscape. They do not uniquely identify
the self-model mechanism, demonstrate calibrated low margins for a majority,
or validate the Gaussian assumption below. A direct assessment would measure
choices and available alternatives under matched interventions, distinguish
changes in information from changes in classification, and estimate the
margin distribution instead of presupposing its shape. Habituation requires a
changing-state model; the static common-shift calculation is one conditional
slice of that wider process.

\subsection{A population model of ambivalence and shifted thresholds}

The proposed population hypothesis is that most people are ambivalent in their
application of ethical constraints, rather than committed to destructive conduct.
In the present geometry, a change in classification can therefore move ordinary
actors across a decision threshold by lowering represented loss. This hypothesis
concerns the population's response to the account. The derivation of occurring
consequences does not depend on how the population is distributed.

For a calculable conditional model, fix one harmful action $h$ and a comparison
action $b$ under stated conditions. For actor $j$, define the converted decision
margin
\begin{equation}\label{eq:population-margin}
m_j=\Phi_{M_j}(h\mid x_j)-\Phi_{M_j}(b\mid x_j).
\end{equation}
Positive $m_j$ favours the comparison action under the assumed minimising policy;
negative $m_j$ favours $h$. Use a common calibrated score unit and first let the
margin distribution have any continuous cumulative distribution function $F$.
A classification intervention reducing margins by a common $\delta\geq0$ gives
$m'_j=m_j-\delta$ and hence
\begin{equation}\label{eq:population-general}
\Pr(m'_j<0)=F(\delta).
\end{equation}
The newly crossing fraction is $F(\delta)-F(0)$. Continuity removes tie mass;
a distribution with atoms requires an explicit tie rule and the corresponding
one-sided probabilities. The threshold-shift mechanism therefore does not
require a normal distribution.

For the proposed normal case, assume $m_j\sim\mathcal N(\mu,\sigma^2)$,
$\sigma>0$, and write $F_{\mathrm N}$ for the standard normal cumulative
distribution function. This is a distribution of conditional decision margins,
not a numerical measure of goodness or evil. Normality remains a modelling
assumption rather than an empirical result established by the cited studies.
The conditional proportion selecting $h$ is
\begin{equation}\label{eq:population-response}
p_h(\delta)=\Pr(m_j<\delta)
=F_{\mathrm N}\!\left(\frac{\delta-\mu}{\sigma}\right).
\end{equation}
Thus $\mu>0$ gives a baseline in which fewer than half select $h$, while a shift
$\delta>\mu$ makes the same action favoured by more than half in the conditional
model. The actors need not acquire a new commitment to cruelty: their represented
comparison has changed. When $\mu>0$, the admissible shift $\delta=\mu$
gives a predicted fraction of one half, and sensitivity to an additional common
shift is greatest there. The calculation
holds the available actions and other margin components fixed; coercion,
heterogeneous shifts, changed opportunities, or social feedback require those
features to be modelled separately. Only actors with $0<m_j<\delta$ newly
cross the threshold in this direction; their fraction is
$p_h(\delta)-p_h(0)$. Actors already favouring $h$ contribute to the resulting
fraction but have not been recruited by this shift.

To make ``ambivalent majority'' inspectable, declare a margin width
$\varepsilon>0$ and use $|m_j|\leq\varepsilon$ as a proxy for choices readily
changed by a shift of comparable size. Its population fraction is
\begin{equation}\label{eq:population-ambivalence}
p_{\mathrm{near}}=
F_{\mathrm N}\!\left(\frac{\varepsilon-\mu}{\sigma}\right)
-F_{\mathrm N}\!\left(\frac{-\varepsilon-\mu}{\sigma}\right).
\end{equation}
The proposed majority condition is $p_{\mathrm{near}}>1/2$. It is a condition on
$\mu$, $\sigma$, and a substantively justified width, rather than a consequence
of normality alone. A small net margin may also arise from strong opposing
commitments; identifying psychological ambivalence requires further evidence.
These quantities describe possible participation under matched conditions,
without inferring universal criminal capability or intrinsic moral character.
The reference function \texttt{normal\_response} calculates the conditional
fractions from declared parameters; it does not fit or validate them.

Granovetter's threshold models are direct prior art for aggregate changes arising
from distributions of individual decision thresholds
\citep{granovetter1978threshold}. His actors' thresholds depend on others'
participation, whereas the calculation here holds the common shift exogenous.
The informational contribution proposed here connects a classification-induced
shift to the losses omitted from the actor's model, while retaining those losses
in the whole-system account. A majority participating in the resulting action
would not, by itself, establish either the assumed normal distribution or a
majority that originally desired its consequences. In this mechanism, actors
who actively deform the field can change the choices of a broad population
with small margins, without each influenced actor acquiring an independent
commitment to the destructive outcome.

\subsection{Three different quantities}

First, there is the consequence account of the proposed action: which continuing
lives, relations, capabilities, and resources it changes. Second, there is the operative
valuation: how the self actually weighs or orders losses and burdens. Its
self-model can admit, omit, or reject a representation of that valuation. Third,
there is the consequence account of requiring the representation to become
explicit. The life affected by the action, the number written by the actor,
and the actor changed by writing it occupy different places in the account.

A number can express a grounded quantity, a decision conversion, or an estimate
with uncertainty. The label ``the value of a human life'' does not specify which
of these it represents. The present framework derives an affected person's
consequence account under a contract; a monetary conversion of that account is
an additional declared evaluation. This distinction gives the written number
an inspectable meaning without making a universal scalar an untyped property
of every life.

A self-model's declaration of pricelessness leaves several possibilities open:
a genuine decision constraint, an unresolved conversion, a socially maintained
commitment, or a compression used alongside finite practical choices. A declaration
alone does not identify which one holds. Neither does every finite resource
choice identify a stable hidden price: the inference requires a decision model.

\subsection{Quantifying an implied conversion}

For a simple conditional model, let an intervention avoid $d>0$ units of a
declared loss measure and require $c\geq0$ additional resource units. Let its
score difference from the baseline be $c-\lambda d$, where $\lambda\geq0$ converts
loss units into those resource units. With other effects held fixed, accepting
the intervention implies $\lambda\geq c/d$ and rejecting it implies
$\lambda\leq c/d$, allowing equality at a tie. Several choices therefore imply
\begin{equation}\label{eq:shadow}
\max_{\mathrm{accepted}}c/d\ \leq\lambda\leq\
\min_{\mathrm{rejected}}c/d.
\end{equation}
An empty accepted set gives no positive lower bound; an empty rejected set
gives no finite upper bound. Inconsistent bounds reject this fixed scalar
model. The inference concerns the declared loss unit and these choices; it is
not a recovered intrinsic price of a particular person.

For example, acceptance at three resource units per avoided loss unit and
rejection at five give $3\leq\lambda\leq5$. Writing $\lambda=4$ proposes a specific
conversion compatible with those choices. It does not establish that four is
the operative value; the choices identify an interval. A self-model can reject
every finite member of the interval while the operative choices still delimit it. The writing does not alter the
original recipient's already fixed consequences. It can alter the writer's
organisation, their future choices, and the relations in which the written
comparison is available. A policy constraint or nonlinear rule need not admit
the scalar inference at all. The reference function \texttt{shadow\_price\_bounds}
implements the finite interval calculation, including inconsistent inputs.

\subsection{Quantifying what the self-model will not admit}

Choose a declared finite grid $H$ of candidate conversions. Let
$\Lambda_{\mathrm{op}}\subseteq H$ contain the values compatible with the operative
choices under the stated decision model. Let $A_M\subseteq H$ contain the values
the self-model explicitly admits as descriptions of that same operative
valuation, under matched conditions. Asking whether a number describes the
self's choices differs from asking whether the person endorses it as an ideal.

For nonempty $\Lambda_{\mathrm{op}}$, define
\begin{align}
E_M&=\Lambda_{\mathrm{op}}\setminus A_M,\label{eq:model-exclusion}\\
g_M&=\mathbf{1}\{\Lambda_{\mathrm{op}}\cap A_M=\varnothing\},\label{eq:model-gap}\\
f_M&=\frac{|E_M|}{|\Lambda_{\mathrm{op}}|}.\label{eq:model-fraction}
\end{align}
The excluded set preserves which candidates are rejected. The indicator records
complete incompatibility within the declared grid; the fraction is a uniform
candidate-count summary of rejection. It is not a measure of discomfort's
intensity, nor a moral weighting of candidate values. Refining or resampling
the grid changes this summary's contract.

For $H=\{3,4,5\}$ and operative bounds $3\leq\lambda\leq5$, rejection of all three
as possible descriptions gives $E_M=H$, $g_M=1$, and $f_M=1$. The operation need
not introduce new external facts. It forces the self-model to confront the
valuation its own choices already constrain. If no candidate matches the choices,
the grid or decision model has failed; the gap is undefined, not evidence that
the self-model rejects a recovered valuation. If some candidates have not been
probed, their admission status remains unknown. The reference function
\texttt{valuation\_gap} implements the finite comparison for fully classified
candidates.

This is an operational quantification of the proposed distinction. It locates
a mismatch between operative valuation and admitted self-description without
identifying a unique brain variable. Explaining the mismatch specifically by
discomfort requires independent evidence of aversion or avoidance, with audience,
incentives, misunderstanding, and changed decision conditions distinguished.
The cost of sustaining or resolving the mismatch is a further consequence
account, rather than another name for the indicator.

\subsection{Quantifying confrontation and acknowledgement}

Let $e_0$ be the condition in which operative valuation remains unadmitted by
the self-model and $e_v$ the condition in which the actor confronts or writes
an explicit candidate number $v$. Hold the compared external options, resource facts, and outcome beliefs
fixed where the experiment permits. Apply the existing functions to
\begin{equation}\label{eq:numeric-account}
\C^{\mathrm{number}}_\theta(v)=\C_\theta(e_v\mid e_0).
\end{equation}
The condition $e_v$ means confrontation or expression, not demonstrated
admission. Someone can write a number while still rejecting it as a description
of themselves. Acknowledgement requires a further observed change in $A_M$ or
other grounded evidence of changed self-modelling. For a verified admission
condition $e_{\mathrm{admit}}$, take the additional comparison
$\C_\theta(e_{\mathrm{admit}}\mid e_0)$; writing alone does not establish it.

This account records changes in the actor's and other affected frames' futures,
access prices, expenditure, and course profiles. It can contain gained capability
through clearer planning as well as lost trust or increased burden. No assumption
of injury is built into the comparison.

Private reading, private writing, and publication are distinct conditions.
Publication adds audiences and institutional consequences; requiring authorship
can add responsibility or commitment. If the explicit condition also supplies
new facts, changes incentives, or changes the audience, its account includes
those effects. It cannot then be described as a pure effect of numerical format.

As an additional behavioural dimension, for aligned action distributions under
matched circumstances $x$, define
\begin{equation}\label{eq:policy-distance}
B(v;x)=\tfrac12\sum_u|\pi(u\mid x,e_v)-\pi(u\mid x,e_0)|.
\end{equation}
This total-variation distance lies between zero and one and measures changed
action probabilities. It is neither a moral weight nor a universal depth-of-knowledge
score. Equal action distributions can coexist with changed self-relations,
expenditure, or effects in another context. The function \texttt{policy\_distance}
computes a floating-point approximation with normalisation tolerance
$10^{-12}$; the exact bound applies to exactly normalised distributions.
The existing functions compute
Eq.~\eqref{eq:numeric-account} from the grounded branch models.

In this example, ``knowing'' versus ``\emph{knowing}'' distinguishes carrying an
operative valuation from admitting it into the self-model. The gap can already
exist before a number is written. Acknowledgement changes the modelling; because
the self-model is part of the self, that change can also alter later action,
self-relations, and expenditure. It need not change the historical consequences
whose valuation was previously denied.

Moral injury is one possible consequence when confronting or revising this
self-description damages enabling self-relations or trust. Acknowledgement can
also improve calibration and make hidden costs accountable. The operative
valuation, its rejection, and the cost of admitting or maintaining the rejection
are separate objects within the whole-system account.

\section{Compression and ethical behaviour}\label{sec:compression}

A calculable principle need not be explicitly calculated whenever a system acts.
Organisms coordinate, avoid damage, and help others before acquiring a theory of
ethics. The informational account proposes an underlying structure that can be
compressed into affect, habits, expectations, and social practices.

\begin{hypothesis}[Ethical behaviour as compression]\label{hyp:compression}
A substantial part of human and animal ethical behaviour expresses compressed
sensitivity to changes in affected capabilities, enabling relations, and burdens.
The expression depends on available information, learned compression, and
relative decision weightings. It need not represent a formal consequence account
or reproduce the functions printed in this paper.
\end{hypothesis}

Four objects must remain distinct: the actual consequence account $\C$, the
information available to the actor, a compressed representation of that
information, and the actor's decision weights. A poor report can hide a real
cost. A coarse compression can merge different costs. Different weights can
produce different responses to the same represented costs. These are different
explanations and admit different corrections.

\subsection{What compression can preserve}

Let $z$ be a finite consequence feature vector and let $Q(z)$ be a compressed
record. For a fixed decision map $d$, exact preservation requires the decision
to be recoverable from that record.

\begin{proposition}[Exact decision preservation]\label{prop:compression}
There exists a map $g$ with $d=g\circ Q$ if and only if $d$ is constant on every
fibre of $Q$: whenever $Q(z)=Q(z')$, then $d(z)=d(z')$.
\end{proposition}
\begin{proof}
If $d=g\circ Q$, equal compressed records give equal decisions. Conversely,
if $d$ is constant on each fibre, define $g(q)$ as the common decision of the
vectors mapping to $q$. The condition makes this definition well-defined on
the image of $Q$.
\end{proof}

For a concrete finite example, suppose a system compares actions using
Eq.~\eqref{eq:example} with weights $(1,1,2)$. Compression to the single feature
$Q(z)=L_A+K_{B,r}-2G_B$ preserves its score and all score-based choices, although
it discards the separate effects. That compression does not generally preserve
choices after the weights change. Preserving a decision is weaker than preserving
the full cost account.

Approximate compression also has a precise boundary. Let $\widehat z(a)$ be a
reconstructed feature vector and suppose
$|w\cdot\widehat z(a)-w\cdot z(a)|\leq\epsilon$ for every feasible action.
If $\widehat a$ minimises the reconstructed score and $a^*$ the actual score,
then
\begin{equation}\label{eq:regret}
w\cdot z(\widehat a)-w\cdot z(a^*)\leq2\epsilon.
\end{equation}
Indeed, the actual score of $\widehat a$ is at most its reconstructed score plus
$\epsilon$; minimisation bounds that reconstructed score by the one at $a^*$;
and reconstruction contributes at most a further $\epsilon$. If component errors
satisfy $|\widehat z_j-z_j|\leq e_j$, one admissible bound is
$\epsilon=\sum_j|w_j|e_j$.

These results demonstrate how a principle can be expressed through compressed
records, and what such expression loses. They do not establish that brains use
this particular feature vector, compressor, or minimisation rule. The hypothesis
concerns the informational principle expressed through behaviour, with neural
implementation left open.

\subsection{Evidence for ethical behaviour without explicit theory}

Haidt's social intuitionist account places moral judgement partly before explicit
reasoning \citep{haidt2001intuition}. Instrumental helping in young children and
chimpanzees \citep{warnekentomasello2006helping}, and rats avoiding an action that
harms a neighbour \citep{hernandezlallement2020harm}, show behaviours that do not
require an articulated ethical theory. Their evidential role is to support the
possibility of ethical action preceding ethical explanation.

The informational hypothesis adds a proposed organising principle for such
behaviour. Helping can restore another's action possibilities; avoidance can
prevent an imposed burden; reciprocity can preserve enabling relations. The
studies do not directly measure the full continuation account. Their behaviour
can also be explained through learning, affect, affiliation, and reward
expectation. Distinguishing explanations requires predictions about changed
actual futures, changed information, and changed weights.

Computational moral-choice research already fits relative sensitivities to
outcomes. Studies of harms to oneself and others
\citep{crockett2014harm} and of different moral strategies
\citep{vanbaar2019moralstrategies} demonstrate implementable decision models.
The contribution here is not the invention of weighted moral computation. It
is a cost account derived before those weights, together with a compression
interpretation of the resulting behaviour.

Animal paradigms also require care about which variable changes. Reward
attainability can explain responses attributed to inequity in dogs
\citep{mcgetrick2023attainability}. This is relevant to the present framework:
an experiment that changes attainable outcomes must distinguish that change
from a changed comparison weight. Behaviour alone does not identify which
account or priority an animal uses.

\subsection{Graded treatment of entities}

Human moral consideration often varies with perceived organisation, agency,
experience, persistent relationships, and similarity. Mind-perception research
relates attribution of agency and experience to different aspects of moral
judgement \citep{gray2007mindperception}. Empathy and compassion across species
also vary with evolutionary divergence \citep{miralles2019empathy}. These findings
are compatible with coarse proxies for informational organisation, but they do
not directly measure information-defined cost.

The hypothesis explains such proxies as compression: visible similarity,
familiarity, responsiveness, and continuity can make another's future structure
easier to represent. The proxy can be informative, biased, or misleading. A
familiar agent with few affected continuations and an unfamiliar agent with many
can reverse the ranking suggested by appearance. The account remains tied to
the actual relations while the actor's representation can fail to track them.

Food choices illustrate the distinction. People can give different priorities
to ending an animal's continuing organisation, producing milk or eggs, harming
plants, and preserving their own nutrition or cultural practices. Vegetarianism
and veganism therefore need not express identical weights or identical causal
accounts. A YouGov survey reported more self-identified vegetarians than vegans
in Great Britain and the United States, with different margins
\citep{nguyen2024dietarytrends}. This is contextual evidence, not a universal
ordering of diets or a measurement of the animals' informational cost.

The proposed trend is greater consideration for entities represented as bearing
richer consequential organisation and relationships. It is not a rule that
humans consistently maximise that quantity, a claim that food practices are
uniform, or a hierarchy of reality. Its empirical test must separate perceived
structure, actual affected futures, and competing weights. The game case makes
that separation unusually inspectable because the underlying continuation model
can be available.

\section{Relation to earlier work}\label{sec:related}

The ingredients have substantial predecessors. The comparison concerns what
each theory takes as given and what it calculates, rather than a claim that
information has only now entered ethics.

\subsection{Information ethics and cybernetics}

Floridi's information ethics extends ethical consideration beyond a human-centred
population \citep{floridi1999informationethics}. Floridi and Sanders connect evil
to impoverishment of informational structure and state duties concerning entropy
\citep{floridisanders1999entropy}. Their entropy is not simply raw Shannon
uncertainty. The close overlap is an informational, broadly ontological basis
for ethical consideration. The present proposal develops attributed continuation
and expenditure functions while placing priorities in a subsequent decision
rule; it does not import a duty to maximise every entity's informational state.

Von Foerster connects ethics with increasing available choices
\citep{vonfoerster2003ethics}. Ashby's variety supplies an established relation
between available distinctions and regulation \citep{ashby1956cybernetics}.
The present derivation calculates which options and enabling conditions change.
A choice-increasing imperative would be a further evaluation rule over that
account. Causal entropic forces likewise demonstrate dynamics organised around
future path entropy \citep{wissnergrossfreer2013causal}; the adopted maximisation
principle is not by itself a cross-agent cost allocation.

\subsection{Empowerment and reachability-based impact}

Empowerment gives an operational information-theoretic measure of control.
Salge and Polani explicitly use human, robot and transfer empowerment in a
proposal for robotic ethics, including weights and human priority
\citep{salgepolani2017laws}. This is a close and implementable predecessor.
The present account retains attributed losses, gains, expenditures and courses
without imposing a human-serving mission as a condition of its calculation.

Relative reachability compares side effects against a baseline
\citep{krakovna2018reachability}. Attainable utility preservation measures changes
in capacities to optimise auxiliary rewards \citep{turner2020}. Both precede
this paper's use of future capability as an impact measure. Their baseline and
proxy concerns apply here: a convenient baseline can conceal a burden, and a
proxy can conceal substitution. Retaining outcome identity, recipient incidence,
and path effects exposes those omissions; it does not make every finite model
complete.

Kolchinsky and Wolpert derive semantic information through correlations causally
necessary for maintaining a system's existence
\citep{kolchinskywolpert2018semantic}. Their viability criterion connects
informational significance to dynamics across living and nonliving systems.
This paper instead compares affected continuations and enabling expenditures
under declared contracts. Sen's capability approach also centres what an agent
can actually do \citep{sen1980equality}; its evaluative account supplies another
important point of comparison without being the source of the finite set
derivation.

\subsection{Relation to the ethics already present in this series}

The earlier papers already contain parts of the account. Paper 01 locates
expenditure and attributed harm in affected frames. Paper 01b's communication
analysis distinguishes structural influence from goal-relative beneficiary
comparison. Its requirement for represented preferences concerns the latter
comparison; it is not an eligibility condition for the capability losses here.
Paper 02's reversibility preference
(Proposition~\ref{SM-prop:reversibility}) is conditional on valuing the retained
optionality. The current paper preserves that condition as a decision claim.

Paper 02b makes restoration without cancellation of an intervening course
explicit. Paper 04 supplies consequential measure and prohibits an ungrounded
sum across frames. Paper 09 limits an institutional application rather than
claiming a complete ethics; Papers 10 and 10a apply named objectives and
priorities to particular futures. Those priorities can specialise the present
decision layer. They are not thereby derived as universal weights.

Paper 11 assembles these distinctions into a consequence account with explicit
functions, proves its conditional invariances, and adds the biological
compression hypothesis. Its novelty claim rests on this arrangement and the
work it does. Neither information ethics nor implementable moral models begin
with this paper.

\section{An application to AI alignment}\label{sec:alignment}

Alignment is an adjacent application. Instruction-following systems trained
with human feedback use judgements to shape a reward or training objective
\citep{ouyang2022feedback}. Constitutional AI uses stated principles to guide
feedback and revisions \citep{bai2022constitutional}. In these representative
approaches, feedback or principles enter the training criterion. They do not
by themselves supply the independently attributed cost construction developed
here.

A system could instead maintain a consequence account across candidate actions,
estimate uncertainty, and expose the effects before its deployment weights
select a policy. Human priorities, legal constraints, or institutional rules
could then be explicit decision conditions. Changing those conditions would not
rewrite the measured closure of another frame's options.

This application does not characterise all alignment research as preference
fitting: empowerment, relative reachability and attainable utility preservation
already overlap with the present machinery. Nor does a better account solve
causal modelling, boundary choice, deceptive reporting, or conflicting goals.
It offers a basis on which those conflicts and omissions can be made inspectable.
An agent instructed to minimise a coarse reported cost can still exploit what
the report leaves out.

\section{Scope, challenges, and conclusion}\label{sec:conclusion}

The paper's central result is an implementable whole-system account of
informational consequences whose occurrence does not require evaluative
recognition. Its strongest claim is that a bearer
need not first be approved as morally eligible: where the relevant consequence
is instantiated, it is available to the accounting procedure. Its strongest
boundary is that measurement still requires a grounding. The account can be
wrong without the consequence waiting for a correct observer.

Four challenges specify where the theory can fail or require revision.
\begin{enumerate}[leftmargin=*]
\item \textbf{Grounding and completeness.} A validated consequence that the
proposed continuation, resource, and path construction cannot represent would
expose a missing dimension. An outcome quotient that merges it would expose an
inadequate resolution, rather than establish zero cost.
\item \textbf{Attribution and continuity.} A population or resource partition
that loses an affected recipient, falsely identifies a replacement, or charges
a shared expense twice invalidates that account. Nested collective accounts
require explicit overlap treatment.
\item \textbf{Mathematical invariance.} Failure of the stated invariances under
their actual hypotheses would invalidate a structural result. Altering the
causal relations or measurement contract falls outside those hypotheses and
must be reported as a changed comparison.
\item \textbf{Compression adequacy.} The biological hypothesis would lose
explanatory value if changes in represented capabilities, enabling relations
and burdens systematically failed to predict behaviour after competing
information and priorities were controlled. Compression and weight specifications must be fixed across the cases used to
test a prediction. Arbitrarily refitting them for every observation is not an
informative confirmation.
\end{enumerate}

The proposed account of the ethical frame and the biological interpretation
remain open to challenge.
A rival explanation can dispute how closely moral dispositions track these
consequences or identify additional causal organisation the account omits.
Those are claims about coverage and explanation, rather than requirements that
an observer confer value before a restriction or expenditure can occur. Moral
injury sharpens this boundary: external benefit and internal damage can coexist,
and a favourable decision score cannot certify the absence of either.

The derivation supplies separate functions for lost and gained futures,
incremental burdens, loss and gain time profiles, and repair. A deterministic specialisation
connects available distinctions to information capacity. Decisions then apply
relative weights. Compression can preserve or approximate those decisions under
stated conditions, providing a viable bridge to ethical behaviour before explicit
ethical reasoning. Recovering the precise operations of a brain is not required
for that bridge or for the account's use.

A future removed from an affected relation is removed whether it is seen,
reported, or valued. That is the proposed lens shift. Whole-system accounting begins
from what changed and who bears it, including the actor and the relations that
sustain its continuations. Ethics becomes an accounting frame
for consequences that narrower accounts exclude. The ontology supplies their
place in the system; studying ethics exposes omissions and the compressed
responses to them. The value placed on responding is a further operation.

\appendix
\section{Finite reference functions}\label{app:functions}

The following standard-library Python functions implement the finite
specialisation. Graph edges encode admissibility; outcomes encode consequential
identity. A budget-limited graph must carry its remaining budget in states.
The reference functions do not discover the causal model, infer a moral
population, or replace unknown inputs with zero.

For the protection case in Section~\ref{sec:protection}, the calls are:
\begin{lstlisting}
b = {"A": {1, 2, 3, 4, 5, 6}, "B": {"x"}}
a = {"A": {1, 2, 3}, "B": {"x", "y", "z", "u"}}
record = consequence_account(b, a, {}, {("B", "resource"): 1})
z = {"A_loss": record["frames"]["A"]["loss"],
     "B_burden": record["resources"][("B", "resource")],
     "B_gain": -record["frames"]["B"]["gain"]}
w = {"A_loss": 1, "B_burden": 1, "B_gain": 2}
# record: A loss 3; B gain 3; B resource difference 1.
# evaluate(z, w) = -2, as obtained in the text.
\end{lstlisting}
For the operative-valuation example of Section~\ref{sec:numeric}, the finite
calls are:
\begin{lstlisting}
bounds = shadow_price_bounds([(1, 3, True), (1, 5, False)])
# Bounds: lower 3, upper 5, consistent True.
candidates = {3, 4, 5}  # Declared finite grid inside those bounds.
gap = valuation_gap(candidates, set())
# Known rejection of every candidate: incompatible True,
# rejection_fraction 1.0, excluded {3, 4, 5}.
# These numbers measure model compatibility, not discomfort intensity.
\end{lstlisting}
These displayed results follow by direct substitution in the defined functions.
The appendix supplies a reference implementation, not a reported software
benchmark or a biological implementation.

\lstinputlisting{../derivations/11_relational_information_ethics_finite.py}

% Bibliography rendered with plainnat from the source .bib in derivations/.
\begin{thebibliography}{52}
\providecommand{\natexlab}[1]{#1}
\providecommand{\url}[1]{\texttt{#1}}
\expandafter\ifx\csname urlstyle\endcsname\relax
  \providecommand{\doi}[1]{doi: #1}\else
  \providecommand{\doi}{doi: \begingroup \urlstyle{rm}\Url}\fi

\bibitem[Asch(1956)]{asch1956independence}
Solomon~E. Asch.
\newblock Studies of independence and conformity: I. A minority of one against
a unanimous majority.
\newblock \emph{Psychological Monographs: General and Applied}, 70(9):1--70, 1956.
\newblock \doi{10.1037/h0093718}.
\newblock URL \url{https://doi.org/10.1037/h0093718}.

\bibitem[Ashby(1956)]{ashby1956cybernetics}
W.~Ross Ashby.
\newblock \emph{An Introduction to Cybernetics}.
\newblock Chapman and Hall, London, 1956.
\newblock URL \url{https://ashby.info/Ashby-Introduction-to-Cybernetics.pdf}.

\bibitem[Bai et~al.(2022)]{bai2022constitutional}
Yuntao Bai et~al.
\newblock Constitutional {AI}: Harmlessness from {AI} feedback, 2022.
\newblock URL \url{https://arxiv.org/abs/2212.08073}.

\bibitem[Bandura(2002)]{bandura2002disengagement}
Albert Bandura.
\newblock Selective moral disengagement in the exercise of moral agency.
\newblock \emph{Journal of Moral Education}, 31(2):101--119, 2002.
\newblock \doi{10.1080/0305724022014322}.
\newblock URL \url{https://doi.org/10.1080/0305724022014322}.

\bibitem[Cocchini et~al.(2009)Cocchini, Beschin, Cameron, Fotopoulou, and Della Sala]{cocchini2009anosognosia}
Gianna Cocchini, Nicoletta Beschin, Annette Cameron, Aikaterini Fotopoulou, and Sergio Della Sala.
\newblock Anosognosia for motor impairment following left brain damage.
\newblock \emph{Neuropsychology}, 23(2):223--230, 2009.
\newblock \doi{10.1037/a0014266}.
\newblock URL \url{https://doi.org/10.1037/a0014266}.

\bibitem[Crockett et~al.(2014)]{crockett2014harm}
Molly~J. Crockett et~al.
\newblock Harm to others outweighs harm to self in moral decision making.
\newblock \emph{Proceedings of the National Academy of Sciences}, 111\penalty0
  (48):\penalty0 17320--17325, 2014.
\newblock \doi{10.1073/pnas.1408988111}.
\newblock URL \url{https://pubmed.ncbi.nlm.nih.gov/25404350/}.

\bibitem[Dana et~al.(2007)]{dana2007wiggle}
Jason Dana, Roberto~A. Weber, and Jason~Xi Kuang.
\newblock Exploiting moral wiggle room: Experiments demonstrating an illusory
preference for fairness.
\newblock \emph{Economic Theory}, 33(1):67--80, 2007.
\newblock \doi{10.1007/s00199-006-0153-z}.
\newblock URL \url{https://doi.org/10.1007/s00199-006-0153-z}.

\bibitem[Fiske and Tetlock(1997)]{fisketetlock1997taboo}
Alan Page Fiske and Philip~E. Tetlock.
\newblock Taboo trade-offs: Reactions to transactions that transgress the spheres of justice.
\newblock \emph{Political Psychology}, 18(2):255--297, 1997.
\newblock \doi{10.1111/0162-895X.00058}.
\newblock URL \url{https://doi.org/10.1111/0162-895X.00058}.

\bibitem[Floridi and Sanders(1999)]{floridisanders1999entropy}
Luciano Floridi and J.~W. Sanders.
\newblock Entropy as evil in information ethics.
\newblock Etica e Politica, author text hosted by the University of Trieste,
  1999.
\newblock URL \url{https://sites.units.it/etica/1999_2/floridi/index.html}.

\bibitem[Floridi(1999)]{floridi1999informationethics}
Luciano Floridi.
\newblock Information ethics: On the philosophical foundation of computer
  ethics.
\newblock \emph{Ethics and Information Technology}, 1, 1999.
\newblock \doi{10.1023/A:1010018611096}.
\newblock URL \url{https://link.springer.com/article/10.1023/A:1010018611096}.

\bibitem[Granovetter(1978)]{granovetter1978threshold}
Mark Granovetter.
\newblock Threshold models of collective behavior.
\newblock \emph{American Journal of Sociology}, 83(6):1420--1443, 1978.
\newblock \doi{10.1086/226707}.
\newblock URL \url{https://doi.org/10.1086/226707}.

\bibitem[Gray et~al.(2007)Gray, Gray, and Wegner]{gray2007mindperception}
Heather~M. Gray, Kurt Gray, and Daniel~M. Wegner.
\newblock Dimensions of mind perception.
\newblock \emph{Science}, 315\penalty0 (5812):\penalty0 619, 2007.
\newblock \doi{10.1126/science.1134475}.
\newblock URL \url{https://doi.org/10.1126/science.1134475}.

\bibitem[Haidt(2001)]{haidt2001intuition}
Jonathan Haidt.
\newblock The emotional dog and its rational tail: A social intuitionist
  approach to moral judgment.
\newblock \emph{Psychological Review}, 108\penalty0 (4):\penalty0 814--834,
  2001.
\newblock \doi{10.1037/0033-295X.108.4.814}.
\newblock URL \url{https://pubmed.ncbi.nlm.nih.gov/11699120/}.

\bibitem[Hartmann and Vorderer(2010)]{hartmannvorderer2010disengagement}
Tilo Hartmann and Peter Vorderer.
\newblock It's okay to shoot a character: Moral disengagement in violent video
  games.
\newblock \emph{Journal of Communication}, 60\penalty0 (1):\penalty0 94--119,
  2010.
\newblock \doi{10.1111/j.1460-2466.2009.01459.x}.
\newblock URL
  \url{https://research.vu.nl/en/publications/its-okay-to-shoot-a-character-moral-disengagement-in-violent-vide/}.

\bibitem[Hauser et~al.(2021)Hauser, Neumann, Marshall, and Mokros]{hauser2021subcomponents}
Nicole~Claire Hauser, Craig~S. Neumann, Julia Marshall, and Andreas Mokros.
\newblock Rational, emotional, or both? Subcomponents of psychopathy predict opposing moral decisions.
\newblock \emph{Behavioral Sciences \& the Law}, 39(5):541--566, 2021.
\newblock \doi{10.1002/bsl.2547}.
\newblock URL \url{https://doi.org/10.1002/bsl.2547}.

\bibitem[Hernandez-Lallement et~al.(2020)Hernandez-Lallement, Attah, Soyman,
  Pinhal, Gazzola, and Keysers]{hernandezlallement2020harm}
Julen Hernandez-Lallement, Augustine~Triumph Attah, Efe Soyman, Cindy~M.
  Pinhal, Valeria Gazzola, and Christian Keysers.
\newblock Harm to others acts as a negative reinforcer in rats.
\newblock \emph{Current Biology}, 30:\penalty0 949--961.e7, 2020.
\newblock \doi{10.1016/j.cub.2020.01.017}.
\newblock URL
  \url{https://pure.knaw.nl/portal/nl/publications/harm-to-others-acts-as-a-negative-reinforcer-in-rats/}.

\bibitem[Kelman(1973)]{kelman1973violence}
Herbert~C. Kelman.
\newblock Violence without moral restraint: Reflections on the dehumanization
of victims and victimizers.
\newblock \emph{Journal of Social Issues}, 29(4):25--61, 1973.
\newblock \doi{10.1111/j.1540-4560.1973.tb00102.x}.
\newblock URL \url{https://doi.org/10.1111/j.1540-4560.1973.tb00102.x}.

\bibitem[Klyubin et~al.(2005)Klyubin, Polani, and
  Nehaniv]{klyubin2005empowerment}
Alexander~S. Klyubin, Daniel Polani, and Chrystopher~L. Nehaniv.
\newblock Empowerment: A universal agent-centric measure of control.
\newblock In \emph{2005 IEEE Congress on Evolutionary Computation}, pages
  128--135, 2005.
\newblock \doi{10.1109/CEC.2005.1554676}.
\newblock URL \url{https://uhra.herts.ac.uk/id/eprint/282/1/901241.pdf}.

\bibitem[Kolchinsky and Wolpert(2018)]{kolchinskywolpert2018semantic}
Artemy Kolchinsky and David~H. Wolpert.
\newblock Semantic information, autonomous agency, and nonequilibrium
  statistical physics.
\newblock \emph{Interface Focus}, 8:\penalty0 20180041, 2018.
\newblock URL \url{https://arxiv.org/abs/1806.08053}.

\bibitem[Krakauer et~al.(2020)Krakauer, Bertschinger, Olbrich, Flack, and
  Ay]{krakauer2020individuality}
David Krakauer, Nils Bertschinger, Eckehard Olbrich, Jessica~C. Flack, and
  Nihat Ay.
\newblock The information theory of individuality.
\newblock \emph{Theory in Biosciences}, 139:\penalty0 209--223, 2020.
\newblock \doi{10.1007/s12064-020-00313-7}.
\newblock URL \url{https://doi.org/10.1007/s12064-020-00313-7}.

\bibitem[Krakovna et~al.(2018)Krakovna, Orseau, Kumar, Martic, and
  Legg]{krakovna2018reachability}
Victoria Krakovna, Laurent Orseau, Ramana Kumar, Miljan Martic, and Shane Legg.
\newblock Penalizing side effects using stepwise relative reachability, 2018.
\newblock URL \url{https://arxiv.org/abs/1806.01186}.
\newblock First submitted 2018; revised version 2, 2019.

\bibitem[Kunda(1990)]{kunda1990motivated}
Ziva Kunda.
\newblock The case for motivated reasoning.
\newblock \emph{Psychological Bulletin}, 108(3):480--498, 1990.
\newblock \doi{10.1037/0033-2909.108.3.480}.
\newblock URL \url{https://doi.org/10.1037/0033-2909.108.3.480}.

\bibitem[Landry et~al.(2022)]{landry2022propaganda}
Alexander~P. Landry, Ram~I. Orr, and Kayla Mere.
\newblock Dehumanization and mass violence: A study of mental state language
in Nazi propaganda (1927--1945).
\newblock \emph{PLOS ONE}, 17(11):e0274957, 2022.
\newblock \doi{10.1371/journal.pone.0274957}.
\newblock URL \url{https://pmc.ncbi.nlm.nih.gov/articles/PMC9645591/}.

\bibitem[Litz et~al.(2009)]{litz2009moralinjury}
Brett~T. Litz, Nathan Stein, Eileen Delaney, Leslie Lebowitz, William~P. Nash,
Caroline Silva, and Shira Maguen.
\newblock Moral injury and moral repair in war veterans: A preliminary model
and intervention strategy.
\newblock \emph{Clinical Psychology Review}, 29(8):695--706, 2009.
\newblock \doi{10.1016/j.cpr.2009.07.003}.
\newblock URL \url{https://doi.org/10.1016/j.cpr.2009.07.003}.

\bibitem[Luft(2015)]{luft2015dynamic}
Aliza Luft.
\newblock Toward a dynamic theory of action at the micro level of genocide:
Killing, desistance, and saving in 1994 Rwanda.
\newblock \emph{Sociological Theory}, 33(2):148--172, 2015.
\newblock \doi{10.1177/0735275115587721}.
\newblock URL \url{https://doi.org/10.1177/0735275115587721}.

\bibitem[Luft(2023)]{luft2023career}
Aliza Luft.
\newblock The moral career of the genocide perpetrator: Cognition, emotions,
and dehumanization as a consequence, not a cause, of violence.
\newblock \emph{Sociological Theory}, 41(4):324--351, 2023.
\newblock \doi{10.1177/07352751231203716}.
\newblock URL \url{https://doi.org/10.1177/07352751231203716}.

\bibitem[McGetrick et~al.(2023)McGetrick, Peters, Korath, Feitsch, Siegmann,
  and Range]{mcgetrick2023attainability}
Jim McGetrick, Hugo Peters, Anna D.~J. Korath, Romana Feitsch, Susanne
  Siegmann, and Friederike Range.
\newblock Perceived reward attainability may underlie dogs' responses in
  inequity paradigms.
\newblock \emph{Scientific Reports}, 13:\penalty0 12066, 2023.
\newblock \doi{10.1038/s41598-023-38836-w}.
\newblock URL \url{https://www.nature.com/articles/s41598-023-38836-w}.

\bibitem[Miralles et~al.(2019)Miralles, Raymond, and
  Lecointre]{miralles2019empathy}
Aur{\'e}lien Miralles, Michel Raymond, and Guillaume Lecointre.
\newblock Empathy and compassion toward other species decrease with
  evolutionary divergence time.
\newblock \emph{Scientific Reports}, 9:\penalty0 19555, 2019.
\newblock \doi{10.1038/s41598-019-56006-9}.
\newblock URL \url{https://doi.org/10.1038/s41598-019-56006-9}.

\bibitem[{National Institute of Mental Health}(n.d.)]{nimhrdoc}
{National Institute of Mental Health}.
\newblock About RDoC.
\newblock Institutional research-framework description, n.d.
\newblock URL \url{https://www.nimh.nih.gov/research/research-funded-by-nimh/rdoc/about-rdoc}.
\newblock Accessed 11 October 2026.

\bibitem[Nguyen(2024)]{nguyen2024dietarytrends}
Hoang Nguyen.
\newblock Market insights into people's food choices in the us and great
  britain.
\newblock YouGov, 2024.
\newblock URL \url{https://yougov.com/articles/48355-us-uk-dietary-trends}.
\newblock Published 17 January 2024; survey report, not a peer-reviewed
  article.

\bibitem[Ouyang et~al.(2022)]{ouyang2022feedback}
Long Ouyang et~al.
\newblock Training language models to follow instructions with human feedback,
  2022.
\newblock URL \url{https://arxiv.org/abs/2203.02155}.

\bibitem[Parslow(2026{\natexlab{a}})]{parslow00relation}
Matthew Parslow.
\newblock Is: Relation without time, cause, or rule.
\newblock Paper 00, manuscript in this series, 2026{\natexlab{a}}.

\bibitem[Parslow(2026{\natexlab{b}})]{parslow01equalisation}
Matthew Parslow.
\newblock The law of graphic equalisation---theory: Minimal recursive dynamics,
  collapse, frame-relative identity, and residual.
\newblock Paper 01, manuscript in this series, 2026{\natexlab{b}}.

\bibitem[Parslow(2026{\natexlab{c}})]{parslow02manifolds}
Matthew Parslow.
\newblock Seldon manifolds---theory: The geometry of reachable futures.
\newblock Paper 02, manuscript in this series, 2026{\natexlab{c}}.

\bibitem[Parslow(2026{\natexlab{d}})]{parslow04narrativium}
Matthew Parslow.
\newblock Narrativium---theory: Frame-relative capacity for consequential
  deformation.
\newblock Paper 04, manuscript in this series, 2026{\natexlab{d}}.

\bibitem[Rai et~al.(2017)]{rai2017dehumanization}
Tage~S. Rai, Piercarlo Valdesolo, and Jesse Graham.
\newblock Dehumanization increases instrumental violence, but not moral violence.
\newblock \emph{Proceedings of the National Academy of Sciences},
114(32):8511--8516, 2017.
\newblock \doi{10.1073/pnas.1705238114}.
\newblock URL \url{https://pmc.ncbi.nlm.nih.gov/articles/PMC5559031/}.

\bibitem[Salge and Polani(2017)]{salgepolani2017laws}
Christoph Salge and Daniel Polani.
\newblock Empowerment as replacement for the three laws of robotics.
\newblock \emph{Frontiers in Robotics and AI}, 4:\penalty0 25, 2017.
\newblock \doi{10.3389/frobt.2017.00025}.
\newblock URL
  \url{https://www.frontiersin.org/journals/robotics-and-ai/articles/10.3389/frobt.2017.00025/full}.

\bibitem[Samuelson(1938)]{samuelson1938behaviour}
Paul~A. Samuelson.
\newblock A note on the pure theory of consumer's behaviour.
\newblock \emph{Economica}, 5(17):61--71, 1938.
\newblock \doi{10.2307/2548836}.
\newblock URL \url{https://www.jstor.org/stable/2548836}.

\bibitem[Sen(1980)]{sen1980equality}
Amartya Sen.
\newblock Equality of what?
\newblock In \emph{Tanner Lectures on Human Values}, volume~1. Cambridge
  University Press, 1980.
\newblock URL
  \url{https://tannerlectures.org/wp-content/uploads/2024/07/sen80.pdf}.
\newblock Lecture delivered at Stanford University, 22 May 1979.

\bibitem[Shannon(1948)]{shannon1948communication}
Claude~E. Shannon.
\newblock A mathematical theory of communication.
\newblock \emph{Bell System Technical Journal}, 27:\penalty0 379--423,
  623--656, 1948.
\newblock URL
  \url{https://www.cs.yale.edu/homes/yry/readings/general/shannon1948.pdf}.

\bibitem[Shay(2014)]{shay2014moralinjury}
Jonathan Shay.
\newblock Moral injury.
\newblock \emph{Psychoanalytic Psychology}, 31(2):182--191, 2014.
\newblock \doi{10.1037/a0036090}.
\newblock URL \url{https://doi.org/10.1037/a0036090}.

\bibitem[Tassy et~al.(2013)Tassy, Deruelle, Mancini, Leistedt, and Wicker]{tassy2013choice}
S\'{e}bastien Tassy, Christine Deruelle, Julien Mancini, Samuel Leistedt, and Bruno Wicker.
\newblock High levels of psychopathic traits alters moral choice but not moral judgment.
\newblock \emph{Frontiers in Human Neuroscience}, 7:229, 2013.
\newblock \doi{10.3389/fnhum.2013.00229}.
\newblock URL \url{https://doi.org/10.3389/fnhum.2013.00229}.

\bibitem[Tetlock et~al.(2000)Tetlock, Kristel, Elson, Green, and Lerner]{tetlock2000unthinkable}
Philip~E. Tetlock, Orie~V. Kristel, S.~Beth Elson, Melanie~C. Green, and Jennifer~S. Lerner.
\newblock The psychology of the unthinkable: Taboo trade-offs, forbidden base rates, and heretical counterfactuals.
\newblock \emph{Journal of Personality and Social Psychology}, 78(5):853--870, 2000.
\newblock \doi{10.1037/0022-3514.78.5.853}.
\newblock URL \url{https://doi.org/10.1037/0022-3514.78.5.853}.

\bibitem[Tetlock et~al.(2017)Tetlock, Mellers, and Scoblic]{tetlock2017sacred}
Philip~E. Tetlock, Barbara~A. Mellers, and J.~Peter Scoblic.
\newblock Sacred versus pseudo-sacred values: How people cope with taboo trade-offs.
\newblock \emph{American Economic Review}, 107(5):96--99, 2017.
\newblock \doi{10.1257/aer.p20171110}.
\newblock URL \url{https://doi.org/10.1257/aer.p20171110}.

\bibitem[Thomson(1985)]{thomson1985trolley}
Judith~Jarvis Thomson.
\newblock The trolley problem.
\newblock \emph{The Yale Law Journal}, 94(6):1395--1415, 1985.
\newblock \doi{10.2307/796133}.
\newblock URL \url{https://www.jstor.org/stable/796133}.

\bibitem[Turner et~al.(2020)Turner, Hadfield-Menell, and Tadepalli]{turner2020}
Alexander~Matt Turner, Dylan Hadfield-Menell, and Prasad Tadepalli.
\newblock Conservative agency via attainable utility preservation.
\newblock In \emph{Proceedings of the AAAI/ACM Conference on AI, Ethics, and
  Society}, 2020.
\newblock \doi{10.1145/3375627.3375851}.
\newblock URL \url{https://arxiv.org/abs/1902.09725}.

\bibitem[{United States Senate Committee on Governmental
  Affairs}(1995)]{senate1995activities}
{United States Senate Committee on Governmental Affairs}.
\newblock Activities of the committee on governmental affairs.
\newblock Technical Report S. Rept. 104-27, United States Senate, 1995.
\newblock URL
  \url{https://www.govinfo.gov/content/pkg/CRPT-104srpt27/html/CRPT-104srpt27.htm}.
\newblock Includes reports of the 1993 and 1994 video-game hearings.

\bibitem[van Baar et~al.(2019)van Baar, Chang, and
  Sanfey]{vanbaar2019moralstrategies}
Jeroen~M. van Baar, Luke~J. Chang, and Alan~G. Sanfey.
\newblock The computational and neural substrates of moral strategies in social
  decision-making.
\newblock \emph{Nature Communications}, 10:\penalty0 1483, 2019.
\newblock \doi{10.1038/s41467-019-09161-6}.
\newblock URL \url{https://www.nature.com/articles/s41467-019-09161-6}.

\bibitem[von Foerster(2003)]{vonfoerster2003ethics}
Heinz von Foerster.
\newblock Ethics and second-order cybernetics.
\newblock In \emph{Understanding Understanding: Essays on Cybernetics and
  Cognition}. Springer, 2003.
\newblock URL \url{https://link.springer.com/book/10.1007/b97451}.
\newblock Archived primary text:
  \url{https://pangaro.com/hciiseminar2019/Heinz_von_Foerster-Ethics_and_Second-order_Cybernetics.pdf}.

\bibitem[Vu et~al.(2023)]{vu2023ignorance}
Linh Vu, Ivan Soraperra, Margarita Leib, Jo{\"e}l van~der Weele, and Shaul Shalvi.
\newblock Ignorance by choice: A meta-analytic review of the underlying motives
of willful ignorance and its consequences.
\newblock \emph{Psychological Bulletin}, 149(9--10):611--635, 2023.
\newblock \doi{10.1037/bul0000398}.
\newblock URL \url{https://doi.org/10.1037/bul0000398}.

\bibitem[Warneken and Tomasello(2006)]{warnekentomasello2006helping}
Felix Warneken and Michael Tomasello.
\newblock Altruistic helping in human infants and young chimpanzees.
\newblock \emph{Science}, 311\penalty0 (5765):\penalty0 1301--1303, 2006.
\newblock \doi{10.1126/science.1121448}.
\newblock URL \url{https://pubmed.ncbi.nlm.nih.gov/16513986/}.

\bibitem[Wissner-Gross and Freer(2013)]{wissnergrossfreer2013causal}
A.~D. Wissner-Gross and C.~E. Freer.
\newblock Causal entropic forces.
\newblock \emph{Physical Review Letters}, 110:\penalty0 168702, 2013.
\newblock \doi{10.1103/PhysRevLett.110.168702}.
\newblock URL \url{https://alexwg.org/publications/PhysRevLett_110-168702.pdf}.

\end{thebibliography}
\end{document}
